\chapter{{Current} {state} {of} {research} {on} {\textit{cualquiera}}}\label{sec:2}

In this chapter, I review the current state of research on modern uses of \textit{cualquiera} in different varieties of Spanish. Many empirical and theoretical studies of \textit{cualquiera} are based on European Spanish, as will be shown in this chapter. There are a few descriptive studies of \textit{cualquiera} in Argentinian Spanish (ArgSp) (see \sectref{sec:2.1.4}). However, these studies do not relate the description of \textit{cualquiera} to the theoretical studies discussed in \sectref{sec:2.2} through \sectref{sec:2.5}. I will fill this gap in Chapters \ref{sec:3} and \ref{sec:4} by analysing the variation between Modern European and Modern Argentinian Spanish in theoretical terms. I will also look at the diachrony of \textit{cualquiera} in Chapters \ref{sec:4} and \ref{sec:5}, which will show the syntactic and semantic conditions under which the evaluative meaning can occur in diachrony.


\section[Informal description of the properties of modern \textit{cualquiera}]{{Informal} {description} {of} {syntactic} {and} {semantic} {properties} {of} {modern} {\textit{cualquiera}}}\label{sec:2.1}

In European Spanish (EurSp), \textit{cualquiera} can be used as a pronoun, determiner, noun, and postnominal modifier. As a pronoun and determiner, it has a restricted distribution with respect to verb mood and modality. However, as I will show, this restricted distribution is not present with \textit{cualquiera} as a noun or postnominal modifier. In modern Argentinian Spanish, the word can also be used as a gradable adjective without the verb mood restrictions found in European Spanish. 
As many uses of \textit{cualquiera} are the same in both varieties (Argentinian and European Spanish) -- besides its use as an evaluative noun with the meaning of ‘nonsense’ and as an adjective in the sense of ‘bad/ordinary’ in Argentinian Spanish but not in European Spanish -- I only explicitly mark the language variety Argentinian Spanish, when these special uses are important for the discussion.

\subsection{{Classification} {of} {indefinites}} \label{sec:2.1.1-1} %no section in manuscript/ changed section label

{\tabref{table:functions_map} shows a classification of indefinites across languages that takes into account their syntactic and semantic properties according to \citet{Haspelmath1997}. This classification will help to characterise the uses of \textit{cualquiera} in synchrony and diachrony.
As can be seen, the evaluative meanings of indefinites as described in \sectref{sec:1} for \textit{cualquiera} and other indefinites in Romance languages are missing.} 


\begin{table}
\begin{tabularx}{\textwidth}{llQ}
\lsptoprule
\textbf{Abbr.} & \textbf{Label} & \textbf{Example} \\
\midrule
SK & specific known & Somebody called. Guess who? \\
SU & specific unknown & I heard something, but I couldn’t tell what it was. \\
IR & irrealis & You must try somewhere else. \\
Q & question & Did anybody tell you anything about it? \\
CA & conditional antecedent & If you see anything, tell me immediately. \\
CO & comparative & In Freiburg, the weather is nicer than anywhere else in Germany. \\
DN & direct negation & John didn’t see anybody. \\
AM & anti-morphic & I don’t think that anybody knows the answer. \\
AA & anti-additive & The gravity of such act goes beyond any justification. \\
FC & free choice & Anybody can solve this problem. \\
UFC & universal free choice & John kissed any woman with red hair. \\
GEN & generic & Any dog has four legs. \\
IND & indiscriminative & I do not want to go to bed with just anyone anymore. \\
\lspbottomrule
\end{tabularx}
\caption {Classification system of FCIs according to \citet{Haspelmath1997} adapted from \citet{Aloni2010}}
\label{table:functions_map}
\end{table}

{One important concept to describe the meaning of indefinites like \textit{cualquiera} is genericity. Genericity refers to the way in which language expresses generalisations about kinds or classes of entities, rather than specific instances of those kinds (\cite{Carlson1977, Krifka1995}). This can involve statements about typical properties or behaviors of a group as a whole, rather than about individual members of that group. Generic sentences are not anchored in any specific point in time or specify any particular referents. For example, \textit{Any dog has four legs} is a characterising sentence because it makes a general statement about dogs as a category. In contrast, \textit{Ben's dog is dangerous} is specific.} 

\subsection{{The} {pronoun} {\textit{cualquiera}} {and} {the} {determiner} {\textit{cualquier}}} \label{sec:2.1.1}

\textit{Cualquiera} can be used as an indefinite pronoun in modern Spanish to express generic statements, i.e. statements about every person or every entity, especially when it is used as a subject in declarative sentences (\citealt{Brucart1999, Sanchezlopez1999, Eberenz2000}, DLE, s.v. \textit{cualquiera}):

\ea
\label{ex:key:20} %ex 20
  \gll Esto lo hace \textit{cualquiera}.\\
this it does \textsc{cualquiera}\\
  \glt ‘Anybody does this.’

  (@ 12)
\z

\ea
\label{ex:key:21} %ex 21
  \gll \textit{Cualquiera} lo hubiera hecho mejor.\\
\textsc{cualquiera} it have.\textsc{SBJV} done better\\
  \glt ‘Anybody would have done it better.’

  (@ 13)
\z

The indefinite pronoun \textit{cualquiera} can be followed by a prepositional phrase [\textit{de} DP] ‘of DP’ (\citealt{Sanchezlopez1999}: 1041, DLE, s.v. \textit{cualquiera}) in which case the DP restricts the quantificational domain of \textit{cualquiera}, e.g. in examples \REF{bkm:Ref72400832} and \REF{bkm:Ref533953098}  the quantification applies to a particular set of reasons and translators, respectively (for the quantificational property of \textit{cualquiera}, see \sectref{sec:2.2}):

\ea %ex 22
\label{bkm:Ref72400832}  \gll \textit{cualquiera} de las razones enumeradas\\
\textsc{cualquiera} of the reasons listed\\
  \glt ‘any of the listed reasons’

  (@ 14a)
\z

\ea %ex 23
\label{bkm:Ref533953098} \gll \textit{cualquiera} de esos traductores\\
\textsc{cualquiera} of those translators\\
  \glt ‘any of those translators’

  (@ 14b)
\z

\textit{Cualquiera} can also be used as a (quantificational) determiner\footnote{{In this book, I will use the terminology of \citet{Zamparelli2000} to refer to quantifiers in prenominal position such as} {\textit{some/any/every N}} {and the Spanish counterparts as (quantificational) determiners.}} similar to \textit{cada} ‘each’ in declarative sentences to express generic statements (\citealt{Sanchezlopez1999}: 1034, 1036; DLE, s.v. \textit{cualquiera}; \citealt{Menéndez-Benito2010}; \citealt{Rivero2011}):

\ea
\label{ex:key:24} %ex 24
  \gll \textit{Cualquier} persona tiene derecho a decidir sobre su cuerpo.\\
\textsc{cualquier} person has right to decide about their body\\
  \glt ‘Any person has the right to decide about her/his body.’
  
  (@ 15)
\z

\ea
\label{ex:key:25} %ex 25
 \gll \textit{Cualquier} ciudadano tiene derecho a elegir.\\
\textsc{cualquier} citizen has right to vote\\
  \glt ‘Any citizen has the right to vote.’

  (@ 14c)
\z

The similar interpretation of \textit{cualquiera} and universal quantifiers such as \textit{cada} ‘each’ and \textit{todo} ‘all’ can be illustrated by the adverbial modification of both universal quantifiers and \textit{cualquiera} by \textit{casi} ‘almost’ and other adverbs such as \textit{excepto} ‘except/but’ that restrict the domain of quantification from every possible element to a smaller part (see \citealt{Arregui2006}, among others):

\ea
\label{ex:key:26} %ex 26
  \gll Casi todos saben hablar inglés.\\
almost everyone know speak English\\
  \glt ‘Almost everbody can speak English.’

  (@ 17)
\z

\ea
\label{ex:key:27} %ex 27
  \gll Casi \textit{cualquiera} puede decirte eso.\\
almost \textsc{cualquiera} can tell.you that\\
  \glt ‘Almost anybody can tell you that.’

  (\citealt{Arregui2006}, e.g. 65 a.)
\z

\textit{Cualquiera} can be used as a determiner inside a Determiner Phrase (DP) that specifies a relative clause with a subjunctive as shown in \REF{bkm:Ref512708}. This relative clause has an indifference interpretation, ‘I don’t care about the flaws he has’ (see DLE, s.v. \textit{cualquiera}):

\ea %ex 28
\label{bkm:Ref512708} \gll \textit{Cualquier} defecto que tenga, yo lo voy a querer igual.\\
\textsc{cualquier} defect that has I him will to love anyway\\
  \glt ‘Whatever flaws he/it has, I will love him/it anyway.’

  (\citealt{Rizzosalierno2013}: 24, @ 18)
\z

Subjunctive relative clauses such as the one in \REF{bkm:Ref512708} are interpreted similar to conditional clauses as in \REF{bkm:Ref513122} (see \citealt{Quer2000}):


\ea %ex 29
\label{bkm:Ref513122} \gll Aunque tenga defectos, yo lo voy a querer igual.\\
although have.\textsc{3SG.SBJV} defects I him will to love anyway\\
  \glt ‘Even if he/it has flaws, I will love him/it anyway.’

  (@ 19)
\z

Both types of conditionals require a subjunctive. The subjunctive includes worlds other than the actual world. Thus, the example in \REF{bkm:Ref512708} not only takes into account the actual world where the person under discussion might have only a specific flaw, but all other flaws in all possible worlds that the speaker can think of (see \citealt{Quer2000}, \citealt{Chierchia2006}, \citealt{Becker2014} on the role of the subjunctive). The effect of the subjunctive in the relative clause is that the assertion ‘I will love him’ is pragmatically strengthened, because it is evaluated with respect to worlds in which the person under discussion might have some or all kinds of flaws. The assertion is considered as true without any restriction.

We now turn to \textit{cualquiera} as a pronoun. In this use, it can have different syntactic functions, such as that of a subject pronoun, as shown in \REF{bkm:Ref533761017}, or as a copular predicate inside a conditional clause with a subjunctive, as in \REF{bkm:Ref533761034}.

\ea %ex 30
\label{bkm:Ref533761017} \gll \textit{Cualquiera} sea el defecto que tiene, yo lo voy a querer igual.\\
\textsc{cualquiera} be.\textsc{3SG.SBVJ} the defect that has I them will to love anyway\\
  \glt ‘Whatever may be the flaw she has, I will love him anyway.’

  (@ 20)
\z

\ea %ex 31
\label{bkm:Ref533761034} \gll Sea \textit{cualquiera} el defecto que tiene, yo lo voy a querer igual.\\
be.\textsc{3SG.SBVJ} \textsc{cualquiera} the defect that has I him will to love anyway\\
  \glt ‘Whatever be the flaw he has, I will love him anyway.’

  (@ 21)
\z

It is a very common observation in the literature that free choice indefinites have a restricted distribution (\citealt{Quer2000}, \citealt{Menéndez-Benito2010}, \citealt{Alonso-OvalleMenéndez-Benito2011,am2018}, among many others). They are usually licensed in modal contexts, e.g. with modal and generic verbs or in the subjunctive (see \chapref{sec:1.2}) and ruled out in episodic\footnote{{I define a context as episodic if the sentence at issue involves an event variable e, and e is existentially closed: \textit{Jean is eating fruits.} ${\exists}$e [eating fruits (e) \& Agent (e, Jean)].}} sentences in English and Spanish, which are usually marked by perfective aspect as in \REF{bkm:Ref41139913} (see \citealt{Menéndez-Benito2010}, \citealt{AloniEtAl2010}, \citealt{Dayal2013}, among others):

\ea %ex 32
\label{bkm:Ref41139913}
\ea[*]{ 
    \gll Juan cogió \textit{cualquiera} de las cartas de esta baraja.\\
            Juan take.\textsc{3sg.pst} \textsc{cualquiera} of the cards from this deck\\
    \glt ‘Juan took any of the cards from this deck.’
    }

\ex[*]{John took any of the cards from this deck.}

    \citep[34]{Menéndez-Benito2010}

\z
\z 

However, as has been noted by \citet{Quer2000}, \citet{Rivero2011}, and \citet{Menéndez-Benito2010}, among others, it is possible to find the determiner \textit{cualquiera} in episodic contexts only if the noun phrase is modified by a relative clause (so-called ‘subtrigging’ examples):


\ea %ex 33
\label{bkm:Ref41139369} 
\ea[*]{  \label{bkm:Ref41139369a}
\gll Juan compró \textit{cualquier} revista.\\
Juan bought \textsc{cualquier} magazine\\
    \glt ‘Juan bought any magazine.’
    }
\ex[]{ \label{bkm:Ref41139369b} \gll Juan compró \textit{cualquier} revista que tuviera fotos de actualidad.\\
Juan bought \textsc{cualquier} magazine that had photos of {current events}\\
    \glt ‘Juan bought whatever magazine had photos of current events.’
        \citep[67]{Rivero2011}
        }
\z
\z

\citet{Quer2000} argues that free choice indefinites (FCIs) like \textit{cualquiera} are licensed by relative clauses, as in example \REF{bkm:Ref41139369b}. These clauses create a context where we talk not just about one specific event, but about many possible situations — for example, things that usually happened in the past. In this context, the meaning of \textit{cualquiera} becomes connected to the idea that something was true in many similar situations or possible worlds, which means different versions of how things could have happened. The clauses as in \REF{bkm:Ref41139369b} introduce a modal context, which is linked to the quantificational meaning of habitual or generic sentences in the past, where a generic operator quantifies over possible worlds or situations (see \citealt{Krifka1995}). The sentence in \REF{bkm:Ref41139369b} has a conditional interpretation associated with generic/characterising statements: if there was a magazine with current events, John would buy it, or every time there was a magazine with photos of current events, John would buy that magazine. The typical cases of licensing of FCIs under subtrigging involve modal contexts such as a conditional-like interpretation. To sum up, given its restriction to modal contexts, \textit{cualquiera} as a determiner and as a pronoun is ungrammatical as the object of a present or past tense indicative verb that refers to a specific event or situation as in \REF{bkm:Ref41139369a}.

What has not been investigated in the literature so far is the use of determiner or pronoun \textit{cualquiera} under indicative copular verbs that predicate over individuals or entities. In the following context, the property of being a man is predicated over the person Juan. In this context, only the indefinite article \textit{un} is judged as acceptable by speakers, whereas the determiner \textit{cualquier} and the pronoun \textit{cualquiera} are not, according to the literature \citet{Alonso-OvalleMenéndez-Benito2011} and speakers' judgments in informal elicitation:

\ea
\label{ex:key:34} %ex 34
  \gll Juan es un hombre (de 30 años).\\
  Juan is a man {of} 30 years\\
  \glt ‘Juan is a (30-year-old) man.’
  
  (@ 22)
\z

\ea
%ex 35
\ea[*]{ \label{bkm:Ref513974}
\gll Juan es \textit{cualquier} hombre (de 30 años).\\
Juan is \textsc{cualquier} man (of 30 years)\\
  \glt ‘Juan is any 30-year-old man.’
  }
\ex[*]{ \label{bkm:Ref513974b}
\gll Juan es \textit{cualquiera}\\
Juan is anyone\\
\glt ‘Juan is anyone.’
}
\z
\z 

One possible explanation of the ungrammaticality in \REF{bkm:Ref513974} is that \textit{cualquier{\slash}cual\-quiera} are not used in the modal or generic context required for its licensing as shown in \REF{ex:key:25} through \REF{bkm:Ref533761034}. However, if \textit{cualquiera} is used in a licensing context such as a partitive phrase as in \REF{bkm:Ref57364726} or a subtrigging clause as in \REF{bkm:Ref57364726-b}, the sentences are grammatical:

\ea %ex 36
\ea \label{bkm:Ref57364726} \gll Si tu hospedaje es \textit{cualquiera} de estos hoteles a continuación ¡nos lo vas a agradecer!\\
if your lodging is \textsc{cualquiera} of these hotels at following us it going to thank\\ 
    \glt ‘If you are staying at any of these hotels mentioned below, you will thank us!’

    (@ 23)
\ex \label{bkm:Ref57364726-b} \gll Sabio es \textit{cualquiera} que sepa seguir los movimientos, las fases de la vida para mantener el yīn y yáng en armonía en su interior.\\
wise is \textsc{cualquiera} who know.\textsc{3SG.SBJV} follow the movements the phases of the life to maintain the yīn and yáng in harmony in their interior\\
    \glt ‘Wise is anyone who knows how to follow the movements, the life cycles…’

    (Genís Sol, \textit{Flujo del espíritu humano}, 35)
\z
\z

To sum up, \textit{cualquiera} as a pronoun or determiner is not licensed in episodic or affirmative predicative contexts without an overt modal.

One rescuing context of the determiner or the pronoun \textit{cualquiera} in predicative contexts without a modal is negation, which shows that \textit{cualquiera} is polarity-sensitive (see \citealt{Chierchia2006}, \sectref{sec:2.2.2} and \sectref{sec:2.4} on FCIs under negation):

\ea %ex 37
\label{ex:key:37}
 \gll El Sr. Miguel Ángel no es \textit{cualquier} empleado de CC.OO, es director del gabinete de estudios de CC.OO.\\
the Mr. Miguel Ángel not is \textsc{cualquier} employee of CC.OO is director of.the office of studios of CC.OO.\\
  \glt ‘Mr. Miguel Ángel is not just any worker of CC.OO, he is the director of CC.OO’s research bureau’

  (@ 24)
\z

\ea %ex 38
\label{ex:key:38}
  \gll La curva que sale no es \textit{cualquier} curva. Es una parábola.\\
the curve that come.out \textsc{neg} is \textsc{cualquier} curve is a parabola\\
  \glt ‘The resulting curve is not just any curve. It is a parabola.’

  (@ 25)
\z

To sum up, \textit{cualquiera} as a pronoun or as a determiner is not permitted with indicative mood, which refers to actual situations or property descriptions in the actual world without subtrigging or negation.

\subsection{{Postnominal} {\textit{cualquiera}}{} {and} {\textit{cualquiera}} {as} {a} {noun}} \label{sec:2.1.2}
\largerpage

As will be shown in this section, postnominal \textit{cualquiera} has a status different from that of the pronominal and the determiner \textit{cualquiera}. The first difference is that postnominal \textit{cualquiera} is usually used in indefinite noun phrases, as in \REF{bkm:Ref533761157} (i.e. the structure \textsc{un} N \textit{cualquiera}). In contrast to the determiner and pronoun \textit{cualquiera}, \textsc{un} N \textit{cualquiera} can be used in contexts that refer to actual situations or properties of actual entities/individuals. The question is thus whether \textsc{un} N \textit{cualquiera} still expresses modality in these contexts.

In \REF{bkm:Ref533761157}, it is asserted that Juan took a card. The semantic contribution of \textit{cualquiera} is to express that this card was taken randomly, which is the so-called random choice interpretation (RCI). It is commonly assumed in the literature that this interpretation correlates with the appearance of volitional verbs, which imply some agent that has desires or goals, such as \textit{coger} ‘to take/to pick’ as in \REF{bkm:Ref533761157}. The RCI implies a modal meaning expressing that the agent Juan could have picked any other card (see \citealt{Menéndez-Benito2005}, \citealt{ChoiRomero2008}, \citealt{Alonso-OvalleMenéndez-Benito2011,am2018}, \citealt{Sanchezlopez1999}: 1041; \citealt{Arregui2006}; \citealt{Rivero2011}):

\ea %ex 39
\label{bkm:Ref533761157} \gll Juan cogió una carta \textit{cualquiera}.\\
Juan picked a card \textsc{cualquiera}.\\
  \glt Random Choice (RC): Juan picked a card and he chose it indiscriminately — he could have picked any other.

  \citep[2]{Menéndez-Benito2005}
\z

However, in affirmative contexts such as in \REF{bkm:Ref533761081}–\REF{bkm:Ref533953228}, with indicative copular\footnote{{The copula} {\textit{be}} {has different uses (see \citealt{Higgins1979}, \citealt{Zamparelli2000}, \citealt{Kellert2015}, and references therein). It can be used predicationally such as} {\textit{Today is a common day (for me)}}{, or it can be used identificationally/specificationally} {\textit{Today is the 21st of July 2010}} {or}{ \textit{Today is Monday}}{. The predicative use usually comes with an indefinite article such as} {\textit{Today is a/*the nice day}}{. Another test to tease the functions apart is to look at the type of the subject of the copula. In the predicational use, the subject is an individual type, e.g}{\textit{. Jean is a man}}{. In the identificational use, the subject is expressed by a demonstrative or deictic pronoun, e.g.} {\textit{This is Jean}} {or}{ \textit{It’s Jean.}} {According to this typology of copular verbs, it is easy to see that} \textsc{un} {N}{ \textit{cualquiera}} {in} {\textit{Juan es un hombre cualquiera}} {‘Juan is an ordinary man’ or in} {\textit{Es un hombre cualquiera}} {should be interpreted predicationally, i.e.} \textsc{un} {N}{ \textit{cualquiera}} {is a copular complement that describes some property of Juan. From this point on, I will use the term predicative} {\textit{cualquiera}} {to refer to} \textsc{un} {N} {\textit{cualquiera}} {as a predicative complement of predicative verbs such as} {\textit{ser}} {‘be’ and similar verbs (}{\textit{parecer}} {‘resemble’, etc.).}}  verbs as in  \REF{bkm:Ref533761081} or with perception verbs such as \textit{mirar} ‘to look (at)’ or \textit{ver} ‘to see’ like in \REF{bkm:Ref533953228}, \textit{cualquiera} has the evaluative interpretation ‘ordinary/common/unremarkable’ (see \citealt{Alonso-OvalleMenéndez-Benito2011,am2018}, DLE, s.v. \textit{cualquiera}):

\ea %ex 40
\label{bkm:Ref533761081} \gll Juan es un estudiante \textit{cualquiera}.\\
Juan is a student \textsc{cualquiera}\\
\glt ‘Juan is an unremarkable student.’ (EVAL)\\
  (\citealt{am2018}: 3)
\z


\ea %ex 41
\label{ex:key:41}
  \gll Eran dos mujeres \textit{cualquiera}.\\
were two women \textsc{cualquiera}\\
  \glt ‘They were two ordinary women.’ (EVAL)

  (\citealt{Rizzosalierno2013}: 24, @ 26)
\z


\ea %ex 42
\label{bkm:Ref533953228}  \gll Te miró y vio un hombre \textit{cualquiera}, no vio al diablo.
\\
you looked and saw a man \textsc{cualquiera} \textsc{neg} saw the devil\\
  \glt ‘(S)he looked at you and (s)he saw in you an ordinary man, not the devil.’\footnote{Although \textit{un hombre cualquiera} appears syntactically as a direct object, it can be argued that it functions semantically as a \textit{predicative complement}. In this construction, the NP does not refer to a specific, individuated entity but rather ascribes a \textit{property} or \textit{classification} to the referent of the pronoun \textit{te}. This reading is supported by the contrastive structure of the sentence: \textit{“no vio al diablo”} (‘(s)he didn’t see the devil’), which highlights that the speaker is categorising the addressee not as a specific man, but as \textit{the kind of man who is ordinary or unremarkable}. This interpretation aligns with analyses of \textit{object-predicate structures} found in sentences with verbs of perception or evaluation, such as \textit{Lo consideró un héroe} (‘He considered him a hero’) or \textit{La llamó una tonta} (‘He called her a fool’), where the postverbal NP has a \textit{classificatory} or \textit{attributive} function.} (EVAL)

  (@ 14d)
\z


The evaluative function can also be expressed in a sentence that refers to an episodic context, in which \textit{cualquiera} modifies a noun that must refer to some entity or individual that the speaker knows and evaluates negatively, see the interpretation in \REF{bkm:Ref54686931-1} (note that the native speakers I consulted informally also admit the so-called ignorance interpretation (IGN), as in \REF{bkm:Ref54686931-2}, where the speaker simply does not know the exact nature of the N):

 %ex 43
\ea \label{bkm:Ref54686931}
\gll La carta de ajuste, en la pantalla, se complementaba con la música ramplona de una zarzuela \textit{cualquiera}.\\
the card of adjustment on the screen \textsc{refl} complemented with the music mediocre of a zarzuela \textsc{cualquiera} (@ 14e)\\

\ea \label{bkm:Ref54686931-1} ‘The testcard, in the screen, was complemented with the vulgar sound of some ordinary zarzuela.’   (EVAL)

\ex \label{bkm:Ref54686931-2} ‘…with the vulgar sound of some zarzuela that I don’t know.’ (IGN)

\z
\z 

Another type of evaluative meaning, namely the pejorative or derogatory interpretation, can be observed with \textit{cualquiera} as a noun in \textsc{un} \textit{cualquiera},\footnote{{One could assume that} {\textit{cualquiera}} {in} {\textit{un(a) cualquiera}} {was historically interpreted as an adjective with an elided noun (e.g. ‘person’). “Nominal gapping is permitted with PPs and postposed restricted adjectives in Spanish, as in} {\textit{El Ø rojo}} {‘the red one’” \citep[34]{Bosque2001}. However, in contemporary Spanish, \textit{un(a) cualquiera} has become lexicalised and typically refers to a ‘prostitute’.}} see the examples in \REF{bkm:Ref150778470} and \REF{bkm:Ref150778472}. The morphological plural form \textit{cualesquiera,} which was used in earlier periods of Spanish (see \sectref{sec:2.5}), is no longer accepted in Modern (European) Spanish, as shown in  \REF{bkm:Ref533761216} (DLE, s.v. \textit{cualquiera}), although it is still marginally found in literary language:\footnote{{In this work, I abstract away from singular and plural differences of} {\textit{cualquiera}} {with respect to their interpretation.}}

\ea %ex 44
\label{bkm:Ref150778470} \gll Juan es un \textit{cualquiera}.\\
Juan is a \textsc{cualquiera}\\
  \glt ‘Juan is a nobody.’

  \citep[36]{Martinez2005}
\z


\ea %ex 45
\label{bkm:Ref150778472} \gll María es una \textit{cualquiera}.\\
María is a \textsc{cualquiera}\\
  \glt ‘María is a prostitute.’

  (@ 27)
\z


\ea %46
\label{bkm:Ref533761216} \gll Son unos/unas \textit{cualquiera/*cualesquiera}.\\
are some \textsc{cualquiera}/\textsc{cualesquiera}\\
  \glt ‘They are nobodies/prostitutes.’

  (@ 28/@ 29)
\z

It seems that no modality is involved in contexts where \textsc{un} (N) \textit{cualquiera} has an evaluative meaning, such as under copular verbs. The question is thus what status \textit{cualquiera} has in these predicative contexts with affirmative sentence mood and whether it is still used as an indefinite element there. It is often assumed in the literature that postnominal \textit{cualquiera} acts as a determiner (see \sectref{sec:2.1.3} on data and discussion of this assumption). Note that the postnominal \textit{cualquiera} cannot be analysed as a (degree) adjective in European Spanish, because it cannot be modified by degree adverbs (see \sectref{sec:2.1.4} for Argentinian Spanish):

\ea \label{ex:key:47} %ex 47
  \gll Un hombre (*muy) \textit{cualquiera}.\\
a man {  very} \textsc{cualquiera}\\
  \glt ‘A man (*very) cualquiera.’

  (@ 4)
\z

This conclusion is confirmed by the absence of morphological agreement. Adjectives usually inflect for person, number, and gender in Modern Spanish in contrast to \textit{cualquiera} which is invariable (\citealt{Seco1986}, s.v \textit{cualquiera}, \citealt{Rainer1993}, s.v. \textit{cualquiera}). In its marginal use in written Spanish, the plural is realised on \textit{cual} (originally a realtive pronoun) and not on the ending -\textit{a} in \textit{cualquiera}, as we would expect if \textit{cualquiera} were lexicalised as an adjective:

\ea \label{ex:key:48} %ex 48
  \gll ordinario, ordinaria, ordinarios, ordinarias\\
  ordinary.\textsc{m.sg.} ordinary.\textsc{f.sg}. ordinary.\textsc{m.pl}. ordinary.\textsc{f.pl.}\\
  \glt ‘ordinary’
\z 

\ea \label{ex:key:49} %ex 49
  \gll cualquiera, cual(es)quiera, *cualquiero, *cual(es)quieros, *cual(es)quieras\\
  cualquiera.\textsc{sg}. cualquiera.\textsc{pl}. { cualquiera.\textsc{m.sg}}. { cualquiera.\textsc{m.pl.}} { cualquiera.\textsc{f.pl.}}\\
  \glt ‘any’
\z

These restrictions can be accounted for diachronically as \textit{cualquiera} is a morphological compound of the relative pronoun \textit{cual} and the verb \textit{querer} (see \chapref{sec:5}). This is why no gender or plural agreement on the ending \textit{{}-a} (originally a subjunctive form) is possible in written Spanish (see \chapref{sec:5}).

However, we do find the plural form -\textit{s} on the ending -\textit{a} in some Latin American varieties of Spanish (e.g. \textit{cualquieras que sean….}, see \citealt{Rainer1993}: 298) and according to \citet{cp2009} few cases of spoken European Spanish (e.g. \textit{unas cualquieras}, see \citealt{cp2009}: 1125), which can be interpreted as a loss of transparency of the word \textit{cualquiera} such that the item can no longer be interpreted compositionally as a compound of the relative pronoun \textit{cual} and the verb \textit{querer} in these varieties. As will be shown in \sectref{sec:2.1.4}, \textit{cualquiera} can be analysed as an adjective in these varieties because of degree modification (see \sectref{sec:2.1.4} for details):

\ea \label{ex:key:50} %ex 50
  \gll La peli es (muy/re) \textit{cualquiera}.\\
the movie is (very/really) \textsc{cualquiera}\\
\glt ‘The movie is (very) ordinary/mediocre.’  (‘okay’ in ArgSp/ *EurSp)

  (@ 5 / @ 6)
\z

The variation in Spanish can be interpreted as a sign that the reanalysis of \textit{cualquiera} from a relative clause into a degree predicate is more advanced in Argentinian Spanish than in European Spanish.

\subsection{{Syntactic} {analyses} {of} {postnominal} {\textit{cualquiera}}} \label{sec:2.1.3}

\hspace*{-.6pt}The syntactic and semantic status of postnominal \textit{cualquiera} is still unclear in the literature. \citet{Lapesa2000[1974]} interprets the definiteness restriction of \textit{un/*el hombre cualquiera} as evidence that \textit{cualquiera} is a quantificational determiner like \textit{cualquier} in \textit{cualquier hombre}. The ‘determiner’ \textit{cualquiera} can co-occur with \textit{un} and not with \textit{el} because \textit{un} does not have any semantic content according to Lapesa (2000 [1974]). Determiners that do not have any semantic meaning are analysed as weak determiners in the literature, whereas determiners that contribute semantic content to the overall meaning are called strong determiners (see \citealt{Zamparelli2000}).

The same observation and conclusion about the determiner status of \textit{cualquiera} has been made by \citet{Rivero2011}, who has shown that postnominal \textit{cualquiera} cannot combine with quantified or definite nouns, but only with bare nouns or indefinite nouns (\textit{un, una,} \textit{unos, unas} N-sg/N-pl) or nouns modified by numerals (\textit{dos, tres,} etc.). The following example in \REF{bkm:Ref63458156} is from \citet[72]{Rivero2011}:

%alternative with a second example
%The same observation and conclusion about the determiner status of \textit{cualquiera} has been made by \citet{Rivero2011}, who has shown that postnominal \textit{cualquiera} cannot combine with quantified or definite nouns, but only with bare nouns or indefinite nouns (\textit{un, una,} \textit{unos, unas} N-sg/N-pl) or nouns modified by numerals (\textit{dos, tres,} etc.). The following examples in \REF{bkm:Ref63458156} and \REF{bkm:Ref533762162} are from \citet[72]{Rivero2011}:
\ea \label{bkm:Ref63458156} %ex 51
\ea[*]{
\gll Puedes comprar \{muchos / algunos / todos / los / estos\} libros cualesquiera.\\
    can buy \hphantom{\{}many ~ some ~ all ~ the ~ these books \textsc{cualesquiera}\\
    \glt ‘You can buy \{many/some/all/the/those\} random books.’
    }
\ex[*]{ \label{bkm:Ref63458156-2}
\gll Puedes comprar \{este / ese / aquel / el~~\} libro \textit{cualquiera}.\\
can buy \hphantom{\{}\textsc{dem.prox} ~ \textsc{dem.med} ~ \textsc{dem.dist} ~ the book \textsc{cualquiera}\\
    \glt ‘You can buy \{this, that, the\} random book.’
    }
\z
\z

%\ea \label{bkm:Ref533762162} %ex 52
%\ea
%\gll Compra (otras) \{dos, tres, cuatro, etc.\} revistas cualesquiera.\\
%buy other \{two, three, four, etc.\} magazines \textsc{cualesquiera}\\
%    \glt ‘Buy any \{two, three, four, etc.\} (more) journals’
%\ex \label{bkm:Ref533762162-2} \gll Compra unas revistas cualesquiera.\\
%\glt buy some magazines \textsc{cualesquiera}.\\
%    \glt ‘Buy any journals.’
%\z
%\z

To conclude, it is assumed in the literature on postnominal \textit{cualquiera} that it is a strong (quantificational) determiner because it cannot co-occur with other determiners of the same kind. As will be shown in \sectref{sec:3.1} on corpus data from Modern Spanish, postnominal \textit{cualquiera} can co-occur with other strong determiners. Similar data will be shown on diachronic corpus data in \sectref{sec:5}. I will conclude that postnominal \textit{cualquiera} is not a determiner.

In order to explain the licensing of postnominal \textit{cualquiera} in episodic contexts, \citet{Rivero2011} assumes that postnominal \textit{cualquiera} encodes lexical modality; that is, \textit{cualquiera} licenses itself and thus does not need a licensor. According to \citet{Rivero2011}, the item \textit{quiera} inside \textit{cualquiera} is interpreted as a modal operator that creates a modal context and licenses \textit{cualquiera}. The following structure is a representation of the Logical From (LF), according to which the item \textit{quiera} is interpreted as a modal operator that scopes over the sentence with a perfective verb:

\ea %ex 53
  LF: [\textsubscript{Sentence} [\textsubscript{Modal} {quiera} ] [\textsubscript{Sentence} Juan compró [\textsubscript{NP} [\textsubscript{NumP} una ] revista cual [\textsubscript{Modal} quiera ]]]]
\z

Although \citet{Rivero2011} did not consider predicative contexts in affirmative sentences in which \textsc{un} N \textit{cualquiera} is also lincensed, one could assume that her explanation covers these cases as well:

\ea %ex 54
  LF: [\textsubscript{Sentence} [\textsubscript{Modal} {quiera} ] [\textsubscript{Sentence} Juan es [\textsubscript{NP} [\textsubscript{NumP} un ] hombre cual [\textsubscript{Modal} quiera ]]]]
\z

However, \citet{Rivero2011} is silent about in what sense \textit{cualquiera} expresses modality in unlicensed contexts such as episodic or predicative contexts.

\subsection{{\textit{Cualquiera}} {in} {Argentinian} {Spanish}} \label{sec:2.1.4}

Before turning to the use of \textit{cualquiera} in Argentinian Spanish, it is important to introduce various forms of linguistic variation in order to understand the difference between European and Argentinian Spanish. Linguistic variation refers to the differences in language use among speakers or within a single speaker in different contexts. Diaphasic, diastratic and diamesic variation are important types of linguistic variation (\citealt{ferguson1994dialect, labov1972sociolinguistic, koch1990gesprochene}). Diaphasic variation, also known as register or style variation, refers to the changes in language use depending on the context or situation. Diamesic variation is influenced by factors such as the formality of the situation or language use that varies on a scale from very formal to very informal or from “Distanzsprache” ‘language of distance’ to “Nähesprache” ‘language of immediacy’ in terms of \citet{koch1990gesprochene} and \citet{koch1999}. For example, a person might use formal language in an administrative/formal context (French: \textit{Que désirez-vous manger?} ‘What do you want to eat?’) or informal language in the communication with friends (\textit{vs. Vous voulez manger quoi? } ‘What do you want to eat?’).

{Diastratic variation, also known as sociolectal variation, refers to the differences in language use associated with social groups (\cite{labov1972sociolinguistic, trudgill2000sociolinguistics}). This type of variation is influenced by factors such as socio-economic status, age and gender.} 

{Argentinian Spanish is influenced by language-contact-induced variation in the context of Italian immigration and by a particular sociolect called \textit{Lunfardo}:\footnote{\textit{Lunfardo}
  {originated from waves of massive immigration from Europe between 1881 and 1914.(see \citealt{Devoto2002}).} {{Lunfardo}} {was originally spoken only in the River Plate region, with a vocabulary enriched from several immigrant languages including Italian \citep{Conde2011}. Sociolinguistic properties of Lunfardo still reflect some connection to Italian immigration (\citet{Kellert-francia-2023}}.}
The status of Italian loanwords and Lunfardo in the present use of Argentinian Spanish became part of the Argentinian identity \citep{Kailuweit2016}. Argentinian Spanish is an interesting variety for linguistic research due to the language contact and sociolect influence on the Spanish variety (see \cite{Conde2011, conde2011lunfardo, engels2011italo, kailuweit2016spirale, wurth2019otro, kellert2023linguistic}, \citealt{Kellert-francia-2023}). The massive immigration of over 4 million Italians to Argentina in the late 19th/early 20th century triggered a well-known case of language-contact-induced variation in Spanish (see \cite{engels2011italo}). The majority of Italian loanwords used in Argentinian Spanish today (e.g., \textit{pibe} ‘little kid, boy’, \textit{laburo} ‘job’) belong to Lunfardo that originated in the lower classes of Buenos Aires (\cite{engels2011italo, wurth2019otro}) and which is thought by some to have emerged among criminals and thieves (\cite{moosekian2016lunfardo}). According to \citet{conde2011lunfardo}, Lunfardo became a “variety of Argentinian Spanish”, which nowadays mainly belongs to colloquial language (\citealt{teruggi1974panorama,FontanelladeWeinberg2000,Wunderlich2012}). Lunfardo and Italian loanwords have been studied based on various data types and methodologies such as a combination of historical and literary texts (\cite{engels2011italo}, \cite{kailuweit2016spirale}), interviews on personal perceptions of Lunfardo (\citealt{wurth2019otro}), word recognition tasks (\cite{moosekian2016lunfardo}) and social media communication data (\citealt{kellert2023linguistic, Kellert-francia-2023}). Some researchers tie the use of Lunfardo and Italian loanwords to social factors like wealth and status. \citet{wurth2019otro} has shown in her study that Lunfardo is primarily used by the lower social classes. \citet{moosekian2016lunfardo} has shown in her study that Lunfardo is no longer used by the younger generation. \citet{kellert2023linguistic, Kellert-francia-2023} have shown a considerable use of Lunfardo words and Italian loanwords in social media communication, thereby confirming the observation that this sociolect is part of everyday communication and part of Argentinian identity as argued by \citet{kailuweit2016spirale}.}

A special use of \textit{cualquiera} is found in Argentinian Spanish according to studies of \citet{Conde2011, Rizzosalierno2013}. These studies describe the usage of \textit{cualquiera} as belonging to informal and colloquial contexts as used primarily by the younger population. 
According to \citet{Conde2011} it has a function of an indefinite pronoun with the meaning similar to \textit{cualquier cosa} ‘anything’ and meanings derived thereof in expressions like \textit{decir cualquiera} or \textit{mandar cualquiera} ‘say something out of place, inconvenient, tell a lie’ or in expressions like \textit{estar en cualquiera} ‘to be in a bad situation, be in trouble’ (literally ‘to be in any [place]’; see also \citealt{Company2009} for \textit{estar en cualquiera}). This use is documented for Lunfardo.

\ea \label{ex:key:55} %ex 55
  Lexical entry of \textit{cualquiera} (\citealt{Conde2011}, \textit{Diccionario etimológico del lunfardo}, s.v.)

  \textit{cualquiera}. pron. indef. En las exprs. ss.: DECIR o MANDAR CUALQUIERA: decir algo fuera de lugar, inconveniente o disparatado; inventar una mentira. {\textbar} 2. ESTAR EN CUALQUIERA: pertenecer a ambientes delictivos o relacionados con la drogadicción; estar distraído. (Del esp. \textit{cualquiera}: persona intedeterminada, que pasa a significar aquí “cualquier cosa”.)

‘\textit{cualquiera.} indef. pronoun. In the expressions such as: DECIR o MANDAR CUALQUIERA: say something out of place, inconvenient, or unreflected; to invent a lie. {\textbar} 2. 2. ESTAR EN CUALQUIERA: to belong to criminal environment or to environments related to drug addiction; to be distracted. (From Sp. \textit{cualquiera}: indeterminate person, which comes to mean here “anything”)’
\z

However, this use of \textit{cualquiera} is not restricted to this variety. According to a note by \citet{ditullio2015cualquiera} in the framework of DILE (\textit{Diccionario Latinoamericano de la Lengua Española}), under verbs such as \textit{decir} or \textit{mandar} ‘to say’, \textit{cualquiera} means ‘nonsense/incorrect/false’:

\ea \label{ex:key:56} %ex 56
  \gll A: Me dijiste que venías temprano. B: \textit{Cualquiera}.\\
    {} me told that would.come early {} \textsc{cualquiera}\\
  \glt A: ‘You said that you would come early.’ B: ‘Nonsense’\\
    (\citealt{ditullio2015cualquiera}; DILE, s.v.)
\z

\ea \label{ex:key:57} %ex 57
  \gll mandaste \textit{cualquiera}\\
  said \textsc{cualquiera}\\
  \glt ‘You said nonsense’\\
(\citealt{ditullio2015cualquiera}; DILE, s.v.)
\z

There is a similar pejorative use of the free choice indefinite pronoun \textit{n’importe quoi} (lit. ‘it doesn’t matter what’) under verbs of saying and copular verbs in Modern French, where it has been nominalised (see \citealt{Paillard1997}, \citealt{Pescarini2009}, \sectref{sec:3} for details):

\ea \label{ex:key:58} %ex 58
  \gll {dire}/{c’est} {(du grand)} \textit{n’importe quoi}.\\
  {say}/{it is} {(of-the big)} \textsc{n’importe quoi}\\
  \glt ‘to say/it’s (total) bullshit.’\\
  (@ 30 / @ 31)
\z

In this use, \textit{cualquiera} can be modified by a number of degree modifiers or affixes. According to \citet{Rizzosalierno2013}, all of them refer to extreme or high degrees (see \citealt{Rizzosalierno2013}: 27).


\ea \label{ex:key:59} %ex 59
  \textit{re-}: prefix\footnote{{I disagree with the analysis of} {\textit{re}} {as a prefix as it is always written as a separate word in social media and can be reinforced by repetition} {\textit{re re cualquiera}} {‘very very bad’.}}  that indicates an extreme degree of a gradable predicate
\ea \textit{muy}, \textit{mal}
\ex\textit{demasiado}, \textit{absolutamente}, \textit{totalmente} ‘too, absolutely, totally’
\ex-\textit{ito}, -\textit{azo}: affective suffixes
\ex-\textit{ísimo}: suffix indicating extreme degree, mostly used with feminine nouns
\z 
\z 

\ea \label{ex:key:60} %ex 60
  \gll \textit{Es} \textit{demasiado} \textit{cualquiera}, casi me hace estrellar mi cabeza contra el monitor.\\
  is too \textsc{cualquiera} almost me makes slam my head against the monitor\\
  \glt ‘It’s too \textit{bad/ordinary}, it almost makes me slam my head against the monitor.’\\
  (\citealt{Rizzosalierno2013}: 27; ArgSp)
\z 

\ea \label{ex:key:61} %ex 61
  \gll Es una mina, \textit{re} \textit{cualquiera.}\\
  is a girl very \textsc{cualquiera}\\
  \glt ‘It is a very ordinary girl.’\\
  (\citealt{Rizzosalierno2013}: 26; ArgSp)
\z 

\largerpage[2]
\ea \label{ex:key:62} %ex 62
  \gll Bueno vieja, pero entonces cambiale el título al thread porque {si no} \textit{es} \textit{cualquiera} \textit{mal.}\\
  well {old lady} but then change the title to.the thread because otherwise is \textsc{cualquiera} bad\\
  \glt ‘Well, buddy, but then change the title of the thread, otherwise it’s really \textit{bad}.’\footnote{{If} {\textit{mal}} {has a status of a degree modifier modifying} {\textit{cualquiera}} {then the syntactic order is inverse in} {\textit{cualquiera mal/cualunque mal}} {as compared to} {\textit{re cualquiera}}{, where the degree modifier is placed before} {\textit{cualquiera}} {rather than after it. In the future, I will test the hypothesis as to whether} {\textit{cualquiera}} {can be analysed as a degree modifier itself in} {\textit{cualquiera mal}}.} \\
  (\citealt{Rizzosalierno2013}: 27; ArgSp)
\z\clearpage

\ea \label{ex:key:63} %ex 63
  \gll Si es un ranking de todos los tiempos este que pusiste es \textit{cualquier}{-ita}!!!\\
  if is a ranking of all the times this that put.2SG is \textsc{cualquiera}\textsc{-ita}\\
  \glt ‘If this is a ranking of every time, this one you posted is \textit{bad}!!!’\\
  (\citealt{Rizzosalierno2013}: 27; ArgSp)
\z 

\ea \label{ex:key:64} %ex 64
  \gll \textit{Cualquier}-aza ese video.\\
  \textsc{cualquiera.affective} that video\\
  \glt ‘That video is extremely \textit{bad}.’\\
  (\citealt{Rizzosalierno2013}: 27; ArgSp)
 \z  

\ea \label{ex:key:65} %ex 65
  \gll Es \textit{cualquier}-ísima tu planteo.\\
  is \textsc{cualquiera.elative} your proposal\\
  \glt ‘Your proposal is extremely \textit{bad}.’\\
  (\citealt{Rizzosalierno2013}: 27; ArgSp)
\z 

\ea \label{ex:key:66} %ex 66
  \gll Es \textit{absolutamente} \textit{cualquiera.}\\
  is absolutely \textsc{cualquiera}\\
  \glt ‘It is absolutely \textit{bad}.’\\
  (\citealt{Rizzosalierno2013}: 28; ArgSp)
\z 

\ea \label{ex:key:67} %ex 67
  \gll Es \textit{totalmente} \textit{cualquiera} eso de haber puesto a Roger con la voz de Syd.\\
  is totally \textsc{cualquiera} that of having put to Roger with the voice of Syd\\
  \glt ‘Having put Roger with Syd’s voice is totally \textit{bad}.’\\
  (\citealt{Rizzosalierno2013}: 28; ArgSp)
\z 

\ea \label{ex:key:68} %ex 68
  \gll Eso de la tanga es \textit{completamente} \textit{cualquiera.}\\
  that of the thong is completely \textsc{cualquiera}\\
  \glt ‘That thong thing is completely \textit{bad}.’\\
  (\citealt{Rizzosalierno2013}: 28; ArgSp)
\z

\textit{Cualquiera} can also have an attributive use in a postnominal position as in European Spanish (see \sectref{sec:2.1.3}):

\ea \label{ex:key:69} %ex 69
  \gll Esto es \textit{un} \textit{video} \textit{cualquiera.}\\
  this is a video \textsc{cualquiera}\\
  \glt ‘This is an ordinary video.’\\
  (\citealt{Rizzosalierno2013}: 26; ArgSp)
\z 

\ea \label{ex:key:70} %ex 70
  \gll Acá les dejo \textit{las} \textit{fotos} \textit{cualquiera} de hoy.\\
  here you leave the photos \textsc{cualquiera} of today\\
  \glt ‘Here you have \textit{ordinary} photos of today.’\\
  (\citealt{Rizzosalierno2013}: 26; ArgSp)
\z

According to \citet{Rizzosalierno2013}, \textit{cualquiera} can also be used under transitive verbs and be interpreted in a pejorative or derogatory way in Argentinian Spanish. However, she does not say anything about its different syntactic function in this context:

\ea \label{ex:key:71} %ex 71
  \gll \textit{Me} \textit{contestó} \textit{cualquiera} o ni me respondía.\\
  me answer.PST.3SG \textsc{cualquiera} or not.even me respond.3SG.PST\\
  \glt ‘(S)he either answered me \textit{bullshit} or not even answered at all.’\\
  (\citealt{Rizzosalierno2013}: 26; ArgSp)
\z 

\ea \label{ex:key:72} %ex 72
  \gll Se ve.3SG que \textit{me} {dijo} \textit{cualquiera} la mina.\\
  it looks that me told \textsc{cualquiera} the girl\\
  \glt ‘Apparently she told me bullshit, that girl.’\\
  (\citealt{Rizzosalierno2013}: 26; ArgSp)
\z

To sum up, \textit{cualquiera} can be used as an adjective with degree affixes and has an evaluative meaning of ‘ordinary/bad’ in Argentinian Spanish (\citealt{Rizzosalierno2013}). It can also be used under transitive verbs such as \textit{decir/mandar} \textit{cualquiera} or \textit{estar en cualquiera} as a direct object comparable to \textit{cualquier cosa} with a deragotary meaning of ‘bullshit’ (\citealt{Conde2011}, \citealt{Tullio2013}, \citealt{Company2009}). %(\tabref{}) 


\begin{table}
\fittable{
\begin{tabular}{lll}
\lsptoprule
 &   {Copular predicate}              &   {Transitive verbs}\\
    \textit{es cualquiera}            &   \textit{decir cualquiera}\\
    \textit{is} \textsc{cualquiera}   &    \textit{to say} \textsc{cualquiera}\\
\midrule
{Evaluative meaning} & + (EVAL1 ‘ordinary, bad’) & + (EVAL2 ‘bullshit’)\\
 {Degree modification} & + & ?\\
\lspbottomrule
\end{tabular}
}
\caption{\label{tab:key:1} \textit{Cualquiera} with evaluative meanings in Argentinian Spanish}
\end{table}  %Table 2

Previous studies on special uses of \textit{cualquiera} in Argentinian Spanish reviewed in this section do not provide a synchronic or diachronic corpus analysis with a quantitative description of various uses and meanings of \textit{cualquiera} in Argentinian Spanish, nor do they discuss the reasons for the synchronic and diachronic variation of \textit{cualquiera} (\tabref{tab:key:1}). The present study aims at filling this gap in the quantitative corpus description of the synchronic and diachronic variation of \textit{cualquiera} in European and Argentinian Spanish by focusing on its use in the synchronic and the historical corpus (see \ref{sec:3} for European Spanish and \ref{sec:4} for Argentinian Spanish). In \chapref{sec:3} and \chapref{sec:4}, I will present new corpus data with a quantitative description, offer a semantic analysis of \textit{cualquiera} in their evaluative uses and discuss reasons for the syntactic and semantic variation of \textit{cualquiera} in European and Argentinian Spanish.

\subsection{{Summary}}  \label{sec:2.1.5}

The table (\tabref{tab:key:2}) sums up the current state of research with respect to different syntactic functions of \textit{cualquiera} and their interpretations. The main observation is that the indicative mood is correlated with different interpretations — EVAL, RC, and ignorance — in both European and Argentinian Spanish. The difference between the two varieties is that in Argentinian Spanish, \textit{cualquiera} with evaluative interpretation is used less restrictively across different syntactic categories, whereas in European Spanish it is only \textsc{un} (N) \textit{cualquiera} that allows the evaluative reading (\tabref{tab:key:1}). 


\begin{table}
\small
\begin{tabularx}{\textwidth}{Q QQ QQ@{}Q@{}Q@{}}
\lsptoprule
 & \multicolumn{2}{c}{ {Syntactic} {context}} & \multicolumn{4}{c}{ {Semantic} {interpretation}}\\
 \cmidrule(lr){2-3}\cmidrule(lr){4-7}
 & {Degree} {modification} & {Indicative} {mood} {in} {episodic} {or} {predicative} {context} & {Evaluative} {interpretation} {(EVAL)} & {Free} {choice} {(FC)} & { {Random} {choice}}\newline {(RC)} & {Ignorance} {interpretation}\\
 \midrule
{Pronoun} & { no} & { no (EurSp)}\newline yes (ArgSp) & { no (EurSp)}\newline yes (ArgSp) & yes & no & no\\
\tablevspace
{\textit{cualquier}} {N} & no & { no (EurSp)}\newline yes (ArgSp) & { no (EurSp)}\newline yes (ArgSp) & yes & no & no\\
\tablevspace
\mbox{\textit{un/una}  (N)}  {\textit{cualquiera}} & { no (EurSp)}\newline yes (ArgSp) & yes & yes & yes (with partitives and under modals) & { yes}\newline{ (episodic tense,}\newline volitional verbs) & yes (episodic tense)\\
\lspbottomrule
\end{tabularx}
\caption{\label{tab:key:2} Syntactic and semantic functions of \textit{cualquiera} in Modern European and Argentinian Spanish}
\end{table} %Table 2

In the following section, I will look more closely at the distribution of different interpretations of \textit{cualquiera} as defined in the literature.

\section{{Formal} {semantic} {analysis} {of} {\textit{cualquiera}}} \label{sec:2.2}

In this section, I will give a brief overview of semantic analysis of \textit{cualquiera} from the literature, which is necessary to understand the formal semantic analysis in this section. I will first introduce the notion of total variation and present formal analyses of \textit{cualquiera}.

\subsection{{Total} {variation} {of} {\textit{cualquiera}} {and} {partial} {variation} {of} {\textit{algún}}}  \label{sec:2.2.1}

\textit{Algún} is an epistemic indefinite that signals the speaker’s ignorance (IGN) with respect to the identity of the individual or object denoted by the indefinite. This makes it incompatible with continuations that specify the referent of the \textit{algún} phrase, as illustrated in \REF{bkm:Ref533761877} (see A\&M 2010 for \textit{algún} phrases, \citealt{Kratzer2002} for German \textit{irgendein}, \citealt{AloniPort2013} for Italian \textit{un qualche,} \citealt{Chierchia2006} for \textit{un} NP \textit{qualunque}, \citealt{JayezTovena2005}, and others for French \textit{quelque} and \textit{un} NP \textit{quelconque}, \citealt{Fălăuş2014} for Romanian \textit{un} NP \textit{oarecare}):

\ea \label{bkm:Ref533761877} %ex 73 
\gll Maria se casó con algún medico. \# En concreto con el Dr. Smith.\\
Maria \textsc{refl} married with some doctor {} in particular with the Dr. Smith\\
  \glt ‘Maria married some doctor. Concretely, Dr. Smith.’
  \textup{(\citealt{am2013}: 36)} \\
\z 

The impossibility of such types of continuation is captured by the “anti-singleton condition”, which expresses the idea captured by \citet{Alonso-OvalleMenéndez-Benito2011} and \citet{GiannakidouQuer2013} that \textit{algún} can only be used if there is more than one alternative known to the speaker (e.g\textit{. algún estudiante} ‘some or other student’). This condition expresses \textit{epistemic indeterminacy} or \textit{anti-specificity} (see also \citealt{GiannakidouQuer2013}, \citealt{Etxeberria2014}). Specificity is given when “a speaker uses an indefinite noun phrase and intends to refer to a particular referent” (\citealt{vonHeusinger2011}: 10). \citet{GiannakidouQuer2013} suggest that the phenomenon of referential vagueness is the absence of specificity, i.e., the absence of referential anchoring and the absence of referential intent. To sum up, \textit{algún} is anti-specific and expresses referential vagueness. It signals the lack of referential intent that characterises specificity (\citealt{vonHeusinger2002,vonHeusinger2007,vonHeusinger2011}). Referential vagueness consists in the minimal variation in possible values, and in uncertainty about which one is the actual value.

One crucial dimension along which epistemic indefinites vary concerns the extent of variation (‘freedom of choice’) imposed on the domain of quantification, which can be \textit{total} or \textit{partial} (see \citealt{Fălăuş2014}, among others, and \sectref{sec:2.2.3}). It has been argued that certain epistemic indefinites (e.g. \textit{irgendein, un} NP \textit{qualunque, un} NP \textit{oarecare}) can sustain TOTAL variation, requiring that \textit{all} relevant alternatives in the domain of quantification qualify as \textit{equally} possible options. For example, the sentence in \REF{bkm:Ref533761905}, with the Romanian epistemic indefinite \textit{un} NP \textit{oarecare}, states that Maria has the obligation to work with a colleague and that this could be any colleague:

\ea  %ex 74
\label{bkm:Ref533761905}  
    \gll Maria \textit{trebuie} să lucreze cu \textit{un} coleg \textbf{\textit{oarecare.}}\\
    Maria must \textsc{sbjv} work with a colleague \textsc{oarecare}\\
    \glt ‘Maria must work with a colleague/any colleague.’\\
    (\citealt{Fălăuş2014}: 32)
\z

In contrast to this total freedom of choice, there are epistemic indefinites that trigger a weaker inference — partial variation — \textit{some, but not necessarily all} alternatives in the relevant domain are epistemic possibilities. As such, they are compatible with the exclusion of some of the possible options. The difference can be most readily observed in the scenario in \REF{bkm:Ref533761939}, as was shown first by \citet{Menéndez-Benito2010} and restated by \citet{Fălăuş2014}:

\ea %ex 75
\label{bkm:Ref533761939}  Maria, Juan, and Pedro are playing hide-and-seek in their country house. Juan is hiding. Pedro believes that Juan is inside the house.

\ea  \gll Juan \textit{tiene} que estar en \textbf{\textit{alguna}} habitación de la casa.\\
Juan has to be in some room of the house\\

\ex \gll Juan \textit{trebuie} să fie în \textbf{\textit{vreo}} cameră din casă.\\
Juan must \textsc{sbjv} be in some room of.the house\\
    \glt ‘Juan must be in a room of the house.’
    (\citealt{Fălăuş2014}: 32)
\z
\z

The context makes it clear that not all rooms of the house are possible choices, that is, the variation associated with the indefinite is limited to a subset of elements in the relevant domain. While the use of \textit{algún} is perfectly acceptable in this scenario, a total variation item like \textit{un} NP \textit{oarecare/qualunque/cualquiera} would be deviant (see \citealt{Fălăuş2014}).

In contrast with \textit{cualquiera}, \textit{algún} requires non-exhaustive variation. The oddness in \REF{bkm:Ref533761963} illustrates a presupposition conflict. The two anti-specific indefinites \textit{algún} and \textit{cualquiera} require opposite things: the former requires non-exhaustive variation,\footnote{Examples expressing non-exhaustive variation in Spanish are expressions like \textit{etcétera} ‘etc.’, \textit{cosas así} ‘this kind of stuff’, \textit{entre otras cosas} ‘among other things’.} while the latter (the free choice indefinite) requires exhaustive variation (but see \sectref{sec:2.3} for grammatical instances of \textit{algún N cualquiera} with evaluative interpretation of \textit{cualquiera}):

\ea %ex 76
\label{bkm:Ref533761963}  
\gll \#Juan ha hablado con algún estudiante \textit{cualquiera}.\\
{Juan} has talked with some student \textsc{cualquiera}\\
\glt ‘{Juan} has spoken with some student.’\\
  (\citealt{Etxeberria2014}: 19)
\z

N \textit{cualquiera} is anti-specific and referentially vague because it is not used referentially and the speaker does not refer to a specific individual (see \citealt{GiannakidouQuer2013}, among others). \citet{Menéndez-Benito2010} assumes that \textit{cualquiera} occurs with nouns that do not denote singletons. The anti-singleton restriction follows from the FCI reading of \textit{cualquiera}, which compares at least two alternatives and conveys that all alternatives count as possible or that the agent does not discriminate between at least two alternatives for his choice. This explanation can explain why \textit{cualquiera} as an FCI does not occur with definite nouns that denote unique individuals like proper nouns do.

Summary of \textit{cualquiera/algún}: both are anti-specific indefinites and both express referential vagueness. They signal the lack of referential intent that characterises specificity. The indefinite \textit{cualquiera} expresses total variation, meaning that all options are evaluated epistemically, whereas in case of \textit{algún} only a subset of possibilities are considered as candidates.

\subsection{{Formal} {analysis} {of} {free} {choice}} \label{sec:2.2.2}
\largerpage
This subsection introduces the formal analysis of Free Choice Indefinites within the Alternative Semantic Framework.

\subsubsection{{Preliminaries: An introduction to alternative semantics}} \label{sec:2.2.2.1}

In order to show how \textit{cualquiera} can be analysed in alternative semantics, I will give a short introduction to the latter based on \citet{Falaus2014}.

On this approach, the semantics of FCIs involves considering alternative semantic values. The theory was developed in a framework of bidimensional semantics (cf. \citealt{Rooth1985}), where alternatives are introduced and computed separately from regular semantic values. The starting point is the interpretation of an exchange such as the following:

\ea \label{ex:key:76} %ex 77
\ea \label{ex:key:76-1} A: Who came to the party?
\ex \label{ex:key:76-2} B: Paul and Sue.
\z
\z

This utterance conveys that Paul and Sue are the only individuals in the context for which a certain property holds (here: having come to the party), a meaning that goes beyond what is explicitly uttered. There are two components leading to this enriched, non-literal interpretation: \REF{ex:key:76-1} a set of alternatives (e.g. contextually relevant individuals) and \REF{ex:key:76-2} the exclusion of the (non-entailed) alternatives. This enriched or strengthened meaning, the one obtained via the consideration and exclusion of alternatives, is often triggered by focus.

Free choice items also activate alternatives. A phrase such as \textit{any student} activates domain alternatives, that is, subsets of (contextually relevant) students (e.g \{Mario\}, \{Marco\}, etc.). Free choice items are indefinites (i.e. existentially quantified elements) that systematically activate sets of alternatives that expand (compositionally) until they find a (quantificational) operator that selects them (e.g. universal quantifier, existential quantifier, negation, question operator) as shown in \REF{bkm:Ref533761448}. Alternatives can be of different types (e.g. individual, propositional) and as a consequence are accessible to both sentential and non-sentential (generalised) quantificational operators, which eventually determine the interpretation of the indefinite. Individual alternatives can expand and give rise to alternatives of a higher type, e.g. propositional alternatives, which then get caught by a sentential operator ${\exists}$p or ${\forall}$p, such as the ones in \REF{bkm:Ref533761448} (see \citealt{Kratzer2002}, \citealt{Chierchia2006}, \citealt{Aloni2007}, \citealt{Alonso-OvalleMenéndez-Benito2011}, among others):

\ea %ex 78
\label{bkm:Ref533761448}  Indefinites in alternative semantics after \citet[8]{Kratzer2002}

  For $⟦α⟧^{w,g} ⊆  D_{<s,t>}:$
\ea{} $⟦α ⟧^{w,g}=\{λw′.${\exists}$p [p${\in}$⟦α⟧^{w,g} \wedge  p(w′) = 1]\}$
\ex{} $⟦α ⟧^{w,g}=\{λw′.${\forall}$p [p${\in}$⟦α⟧^{w,g} →       p(w′) = 1]\}$
\z
\z

The contribution of the indefinite remains constant, namely, to provide a set of alternatives, its different readings being the result of associating with different operators.

When a set of propositional alternatives combines with the existential operator as defined in \REF{bkm:Ref533761448}, it yields the proposition that is true in case at least one alternative in the relevant set is true. In contrast, the universal quantifier in \REF{bkm:Ref533761448} yields the proposition that is true in case every alternative in the relevant set is true.

Now, importantly, as will be shown below in \sectref{sec:2.2.2.2}, \textit{cualquiera} provides alternatives that always combine with the universal operator ${\forall}$p in modal contexts (i.e. under modal verbs, conditionals, subtrigging, and other modalised constructions). Let us see how this approach works for \textit{cualquier libro} ‘some book/any book’ in \REF{bkm:Ref533761676}. The domain of alternatives of \textit{cualquier libro} consists of all individual books of a contextually relevant set, for instance \{\textit{Don Quijote}, \textit{El camino}, …\}. These alternatives grow to become propositions \{You can buy \textit{Don Quijote}, You can buy \textit{El camino}, …\}:

\ea %ex 79
\label{bkm:Ref533761676}  
\gll Puedes tomar \textit{cualquier} libro.\\
can take \textsc{cualquier} book\\
  \glt ‘You can take any book.’

  (@ 32)
\z

In order to derive the free choice meaning in \REF{bkm:Ref533761676} that every alternative (i.e. every book) is a possible option and not just some of them, we compute the enriched meaning of the sentence in \REF{bkm:Ref533761693}, by interpreting the exhaustivity expressed by \textit{cualquier}, i.e. we negate all alternatives that do not match the assertion, which yields the following enriched meaning (summarised under the tenet that all alternatives must be taken into account). The enriched meaning does not contain the statement according to which only one alternative can be taken into account (represented by ${\perp}$ which basically means not-computable):

\ea  %ex 80
\label{bkm:Ref533761693}  Total exhaustivity of alternatives (a ${\vee}$ b ${\vee}$ c) =

    You can take (a ${\vee}$ b ${\vee}$ c) ${\wedge}$

    ¬ You can take (a ${\wedge}$ b) ${\wedge}$

    ¬ You can take (a ${\wedge}$ c) ${\wedge}$

    ¬ You can take (b ${\wedge}$ c) ${\wedge}$

    ¬ You can take (a ${\vee}$ b) ${\wedge}$

    ¬ You can take (b ${\vee}$ c) ${\wedge}$

    ¬ You can take (a ${\vee}$ c) ${\wedge}$

    ¬ You can take a ${\wedge}$ ¬ You can take b ${\wedge}$ ¬ You can take c = ${\perp}$
\z

To sum up, I have discussed so far that \textit{cualquiera} triggers alternatives and total exhaustification over the alternatives in modal contexts (see \sectref{sec:2.2}). In the following section, I will show one possible formalisation of FCIs in Alternative Semantics.

\subsubsection{{\citegen{Aloni2007} formalisation of FCIs}} \label{sec:2.2.2.2}

According to \citet{Aloni2007,AloniPort2013}, and \citet{Aloni2021}, FCIs trigger the application of the total exhaustification over alternatives as demonstrated in \REF{bkm:Ref533761693} and represented formally by an exhaustivity operator {EXH} (see also \citealt{Menéndez-Benito2010}). This operator has a similar function as the focus adverb \textit{only} which associates with alternatives and expresses that only a certain alternative is true. The SHIFT\footnote{Shift means here type shifting, i.e. conversion of one type into another type, e.g. from a type of individuals into a type of propositions.}  function expresses that the exhaustification of domain alternatives shifts to the exhaustification of propositional alternatives, that is, for \REF{bkm:Ref150778689} it is possible that only Peter falls or it is possible that only Mary falls. The total exhaustification expresses that there is no alternative proposition that is left out by the exhaustification:

\ea %ex 81
\label{bkm:Ref150778689}  Anyone can fall. %\citep{AloniEtAl2011}
\ea{}  [${\forall}$](${\lozenge}$(\textsc{SHIFT}‹\textsubscript{s,t}›\textsubscript{} ({EXH}[anyone, fall])))\footnote{the forall-sign in brackets is a typographic convention often used in semantic literature to flag interpretive or pragmatic layers distinct from syntactic structure.}
\ex ${\lozenge}$ only Peter fall  ${\lozenge}$ only Mary fall  ${\lozenge}$ only Bob fall  ...

    \citep{AloniEtAl2011}
\z
\z

She accounts for the distribution of FCIs under modals such as in \REF{bkm:Ref63544535} and negation as in \REF{ex:key:83}, as well as in contexts that also recover universal quantification of exhaustified propositional alternatives without always expressing an overt modal verb such as subtrigging in \REF{bkm:Ref54857683} and degree comparatives as in \REF{bkm:Ref54857674}.

\ea %ex 82
\label{bkm:Ref63544535}   {Canonical FCI}
\ea  Sentence: You can bring any book.
%\ex Logical form: [${\forall}$](${\lozenge}$\textsc{SHIFT}‹\textsubscript{s,t}› ({EXH}[any book, λx. you bring me x]))
\ex Logical form:
[$\forall x$](book(x) → $\lozenge$\textsc{Shift}\textsubscript{⟨s,t⟩}({EXH}[any book, λx. you bring me x]))

\ex Predicted meaning: For each book it is possible that you bring me only that book.
\z
\z


\ea %ex 83
\label{ex:key:83}   {Embedding of the FC effect}
\ea  Sentence: You cannot bring any book.
\ex Logical form: [¬][${\forall}$](${\lozenge}$\textsc{SHIFT}‹\textsubscript{s,t}› ({EXH}[any book, λx. you bring me x]))
\ex Predicted meaning: You cannot freely choose which book you bring me. It is not the case that for each book you can bring me that book.
\z
\z


\ea %ex 84
\label{bkm:Ref54857683}   {Licensing by subtrigging (see \citealt{Aloni2007})}
\ea  Sentence: Anyone who tried to jump fell.
\ex Logical form: [${\forall}$](↓\textsc{SHIFT}\textsubscript{e}({EXH}[anyone, who tried to jump]) fell)
\ex Predicted meaning: All persons who tried to jump fell.
\z
\z


\ea %ex 85
\label{bkm:Ref54857674}   {FCIs in comparatives (see \citealt{Aloni2021})}
\ea  Sentence: John is taller than any girl.
\ex Logical form: [${\forall}$](\textsc{SHIFT}\textsubscript{e}({EXH}[d, λd.tall (j,d)]) > \textsc{SHIFT}\textsubscript{e}({EXH}[d, λd.tall (any girl, d)])
\ex Predicted meaning: For all girls x, John is taller than x.

    \citep{AloniEtAl2011}
\z
\z

\citet{Aloni2007,Aloni2021} and \citet{AloniEtAl2010} also account for the ungrammaticality of FCIs such as \textit{any} in episodic contexts in \REF{bkm:Ref517204} that refer to some factual events by showing that the exhaustivity operator cannot apply there. This is represented as  $\perp$ %⏊
in \REF{bkm:Ref517204}.

\ea %ex 86
\label{bkm:Ref517204}   {Ruling out FCIs in episodic contexts}
\ea  Sentence: \# Anyone fell.
\ex Logical form: [${\forall}$]((\textsc{SHIFT}‹\textsubscript{s,t}› ({EXH}[anyone, fell]))
\ex Predicted meaning: $\perp$%⏊

    \citep{Aloni2021}
\z
\z

Aloni argues that the sentence in \REF{bkm:Ref517204} would give the following incorrect interpretation: for each person only that person fell. A past event described by the perfective aspect can either express that everyone fell or that a certain person/group of people fell. A modal verb as in \textit{Anyone can fall} rescues the interpretation, as it shifts the exaustivity to different worlds (i.e. it becomes interpretable as \textit{For every person it is possible that this person fell}) and not only the actual word (see \citealt{Menéndez-Benito2010}, \citealt{Aloni2021}, among others). This is how these authors account for the relation between modality and free choice indefinites.

Although Aloni makes a correct prediction for FCIs such as English \textit{any} and the determiner/pronoun \textit{cualquiera} in European Spanish, which need a modalised context to be licensed and are not permitted in episodic contexts (see \sectref{sec:2.1}), her analysis does not account for the distribution of (postnominal) \textit{cualquiera} in episodic contexts and with copular verbs in present or past tense in European and Argentinian Spanish (see \sectref{sec:2.1} and \sectref{sec:2.2.1}). The question is how the distribution of \textit{cualquiera} should be accounted for in a system where \textit{cualquiera} triggers alternatives and total exhaustivity over these alternatives. I will deal with the current state research on postnominal \textit{cualquiera} in episodic contexts in \sectref{sec:2.3}.

\subsubsection{{Summary}} \label{sec:2.2.2.3}

It has been shown that, according to the current state of research, FCIs such as \textit{any} or \textit{cualquiera} on their free choice meaning have a universal interpretation if embedded under a possibility modal and other modalised constructions such as subtrigging that recover a possibility modal. According to \citet{Aloni2007}, FCIs trigger total exhaustification over alternatives, captured by an exhaustivity operator EXH. This operator associates with a universal quantifier that operates on propositions. In this way Aloni can account for the universal inference of FCIs, which is a conventional implicature as indicated by the impossibility of cancelling under negation.

\subsection{{Semantic} {analysis} {of} {\textit{cualquiera}} {in} {episodic} {contexts}}  \label{sec:2.2.3}
\subsubsection{{Informal description}} \label{sec:2.2.3.1}

We saw in \sectref{sec:2.2.1} that \textit{cualquiera} is supposed to be excluded from episodic contexts with a free choice interpretation. However, it can express agent-related modality, such as the random choice interpretation. According to \citet{ChoiRomero2008}, \citet{ChoiRomero2008}, and \citet{Alonso-OvalleMenéndez-Benito2011,am2018}, postnominal \textit{cualquiera} used in episodic contexts under volitional verbs has a RC interpretation:

\ea \label{ex:key:87} %ex 87
  \gll Juan necesitaba un pisapapeles, de modo que cogió un libro \textit{cualquiera} de la estantería y lo puso encima de la pila.\\
  Juan needed a paperweight of way that {he took} a book \textsc{cualquiera} from the shelf and it {he put} on.top of the pile\\
  \glt ‘John needed a paperweight, so he took a random book from the shelf and put it on top of the pile.’\\
  (\citealt{ChoiRomero2008}: 97)
\z

\citet{Alonso-OvalleMenéndez-Benito2011,am2018} assume that postnominal \textit{cualquiera} is ambiguous between RC and EVAL (but see \citealt{ChoiRomero2008}, who do not mention the ambiguity):

\ea %ex 88
\label{bkm:Ref533953640} 
\gll Juan compró un libro \textit{cualquiera}.\\
Juan bought a book \textsc{cualquiera}\\
\glt ‘Juan bought a book.’
\ea  John bought a book, and he chose it indiscriminately. (Random Choice)
\ex John bought a book, and the book is not special or remarkable (according to the speaker). (Unremarkable reading, EVAL)

    (\citealt{am2018}: 2)
\z
\z

In order to distinguish EVAL from RC, \citet{Alonso-OvalleMenéndez-Benito2011,am2018} have created different contexts:


\ea %ex 89
\label{bkm:Ref54686706}  Scenario. Juan went to the bookstore. He wanted to buy \textit{The Unbearable Lightness of Being}, and did so. I don’t think this book is special in any way.
\z 

\ea %ex 90
\label{bkm:Ref533953745}  Scenario. Juan went to the bookstore, and bought a book at random. The book turned out to be \textit{The Unbearable Lightness of Being}. I think this book is remarkable.
\z

\begin{quote}
In the scenario in [\REF{bkm:Ref54686706}], I can truthfully utter [\REF{bkm:Ref533953640}] on its unremarkable interpretation (since I consider the book that Juan bought unremarkable) but not on its random choice interpretation (as Juan didn’t make an indiscriminate choice). In contrast, my utterance of [\REF{bkm:Ref54686706}] in the scenario in [\REF{bkm:Ref533953745}] would be false on the unremarkable interpretation but true on the random choice interpretation.
(\citealt{am2018}: 2f.)
\end{quote}

According to \citet{am2018}, on the RC interpretation, the noun modified by \textit{cualquiera} must represent an object that the agent of the sentence must have chosen indiscriminately. They assume that, in this use, \textit{cualquiera} appears in object position of volitional verbs. \citet{Alonso-OvalleMenéndez-Benito2011} assume that the RCI is restricted (e.g. not possible with unaccusative verbs), whereas EVAL is always possible:

\begin{quote}%\todo{check scope}
While the unremarkable interpretation is always possible, the random choice interpretation has a restricted distribution. First of all, for this interpretation to obtain, the sentence has to describe a volitional event (\citealt{ChoiRomero2008}) and \textit{uno cualquiera} should not be in the subject position of agentive verbs. If that condition is not met, \textit{uno cualquiera}\footnote{{A\&M refer to}
  \textsc{un}{} {N}{ \textit{cualquiera}} {as} {\textit{uno cualquiera}}.
 }
only has the unremarkable interpretation. Copular sentences are a clear case. The sentences in (a) and (b) convey that the subject is unremarkable within the class of individuals denoted by the noun phrase: (a) says that Juan is a regular clerk and (b) that Juan is a student just like the rest.
\end{quote}

\ea %ex 91 
\ea  
\gll Juan es un oficinista \textit{cualquiera}.\\
Juan is a clerk \textsc{cualquiera}\\
\glt ‘Juan is an unremarkable clerk.’
\ex 
\gll Juan es un estudiante \textit{cualquiera}.\\
Juan is a student \textsc{cualquiera}\\
\glt ‘Juan is an unremarkable student.’

    (\citealt{am2018}: 3)
\z
\z

The following examples only have the ‘unremarkable’ reading (=EVAL) and not RC, because they do not contain volitional verbs:

\ea %ex 92
\ea \gll Ayer Juan tropezó con un objeto \textit{cualquiera}.\\
    yesterday Juan stumbled with an object \textsc{cualquiera}\\
    \glt ‘Yesterday, Juan stumbled on an unremarkable object.’\\
    (\citealt{am2018}: 3)
\ex \gll La levadura destrozó un molde \textit{cualquiera}.\\ 
    the yeast destroyed a {baking pan} \textsc{cualquiera}\\ 
    \glt ‘The yeast destroyed an unremarkable baking pan.’\\
    (\citealt{am2018}: 3)
\z
\z

Nouns modified by postnominal \textit{cualquiera} that represent subjects of volitional verbs such as \textit{hablar} ‘speak’ or objects of non-volitional verbs cannot have RC as the authors assume, but only EVAL:

\ea %ex 93 
  \gll Habló un estudiante \textit{cualquiera}.\\
  spoke a student \textsc{cualquiera}\\
  \glt ‘An unremarkable student spoke.’ (\# ‘a random student spoke’)
  
  (\citealt{am2018}: 4)
\z

Furthermore, \citet{Alonso-OvalleMenéndez-Benito2011} assume that \textit{cualquiera} can also have a free choice interpretation under some modals, which they call harmonic interpretation (HI):


\begin{quote}
When \textit{uno cualquiera} is embedded under some modals, another possibility arises: \textit{uno cualquiera} introduces a distribution effect with respect to the worlds that the modal ranges over. For instance, \textit{Coge una carta cualquiera!} (‘take any card!’) can be interpreted as conveying that any card is a permitted option.

(\citealt{am2018}: 1)
\end{quote}

\citet{Alonso-OvalleMenéndez-Benito2011} assume that \textit{cualquiera} is lexically ambiguous between RCI/HI on the one hand and EVAL on the other:

\begin{quote}
As a working hypothesis, we will assume that \textit{uno cualquiera} is ambiguous between one form that can give rise to the random choice and harmonic interpretations (the form that we will focus on), and another that conveys the unremarkable interpretation. To determine whether this hypothesis is tenable it would be necessary to provide a detailed analysis of the unremarkable interpretation, a non-trivial task that we will not be able to undertake here.

(\citealt{am2018}: 8)
\end{quote}

To sum up, \citet{Alonso-OvalleMenéndez-Benito2011,am2018} show that RCI is restricted to objects of volitional verbs and it does not arise under non-volitional verbs. The ‘unremarkable’ reading is completely unrestricted; it can appear everywhere:

\begin{table}
\begin{tabularx}{\textwidth}{lQQ}
\lsptoprule
\textbf{Interpretation} & \multicolumn{2}{c}{ \textbf{Restrictions}}\\
\cmidrule(lr){2-3}
& { Available} & Unavailable\\
\midrule
Random choice & {object of volitional verbs} & subject of volitional verbs\newline non-volitional verbs\\
\tablevspace
Unremarkable & unrestricted & unrestricted\\
\lspbottomrule
\end{tabularx}
\caption{\label{tab:key:3} The distribution of RCI and EVAL of \textsc{un} N \textit{cualquiera} (\citealt{Alonso-OvalleMenéndez-Benito2011, AlonsoMenendez2018})}
\end{table} %Table 3

Under their approach, the restriction to volitional verbs follows from the modal component of \textsc{un} N \textit{cualquiera} (cf. \sectref{sec:2.1.3}). The modal component of \textsc{un} N \textit{cualquiera} expresses that every alternative of \textsc{un} N \textit{cualquiera} conforms to the agent’s decision to choose this alternative. Thus, RCI can only apply to volitional verbs, as they express the agent’s decisions to do something (e.g. picking some book).

The random reading disappears whenever there is nothing for the agent to choose, as for example with unaccusative verbs, because unaccusative verbs do not match with the decision condition according to which the agent does something in accordance to his/her decision (see \citealt{Alonso-OvalleMenéndez-Benito2011}, 2018).

\citet{am2018} assume lexical ambiguity between RCI/HI and EVAL and that EVAL is unrestricted. This hypothesis needs to be tested (see \chapref{sec:3}).

\citet[574]{Chierchia2006} observes yet another interpretation of Free Choice indefinites in episodic contexts. He notes that the existential FCI \textsc{un} N \textit{qualsiasi} in Italian can be used in episodic contexts under a covert modal operator that expresses speaker’s modality, namely speaker’s ignorance (see also \sectref{sec:2.2.1} for the ignorance reading expressed by Spanish \textit{algún}):

\ea %ex 94
\label{bkm:Ref533762757} 
\gll Gianni è uscito di corsa e non sapendo che fare, ha bussato ad una porta \textit{qualsiasi}.\\
Gianni is exited in hurry and \textsc{neg} knowing what to.do has knocked at a door \textsc{qualsiasi}\\ 

\glt ‘Gianni ran out and not knowing what to do, knocked on any door.’\\
Speaker’s ignorance: it is possible that for every door x Gianni knocked at the door x.
  \citep[574]{Chierchia2006}
\z

Note that free choice relatives in English can also express an ignorance reading (IGN) in episodic context (see \citealt{Condoravdi2015}: 216), that is, in \REF{ex:key:95} the speaker is ignorant about the identity of the book John bought:

\ea \label{ex:key:95} %ex 95
  Whatever book John bought was very expensive.
\z

According to Chierchia the sentence in \REF{bkm:Ref533762757} could be interpreted as something like “it follows from what the speaker knows that Gianni knocked at a door.” The FCI implicature would then be “it is consistent with what the speaker knows that any door might have been the one knocked at” \citep[574]{Chierchia2006}. Chierchia thus assumes that the speaker considers it a possibility that for every door x, Gianni knocked at the door x. It is important to mention that the modality is associated with the speaker and not with the agent \textit{Gianni}, according to \citet{Chierchia2006} (see \sectref{sec:2.3} for agent’s modality).

Recall from the \sectref{sec:2.1} that native speakers agree on the possibility of an ignorance interpretation of \textsc{un} N \textit{cualquiera} in episodic contexts (see \ref{bkm:Ref54686931}, repeated here in \ref{bkm:Ref72400981}):

\ea %ex 96
\label{bkm:Ref72400981} 
\gll La carta de ajuste, en la pantalla, se complementaba con la música ramplona de una zarzuela \textit{cualquiera}.\\
the test.card of adjustment on the screen \textsc{refl} complemented with the music vulgar of a zarzuela \textsc{cualquiera}\\
\ea ‘The testcard, in the screen, was complemented with the vulgar sound of some ordinary zarzuela.’   (EVAL)

\ex ‘…with the vulgar sound of some zarzuela that I don’t know.’ (IGN)

  (@ 14e)
\z
\z

The ignorance interpretation is not mentioned in the literature on Spanish \textit{cualquiera} (see \citealt{Kellert-francia-2023}, on the ignorance interpretation on a similar item in Argentinian Spanish: \textit{cualunque}). In this book, I do not pursue the study on the ignorance reading in episodic contexts of \textit{cualquiera} either.

\subsubsection{{Formal analysis of the Random Choice interpretation}} \label{sec:2.2.3.2}

The random choice interpretation, RCI, has been described very well in the literature. I will illustrate two prominent proposals, which both assume that postnominal \textit{cualquiera} has a RCI whenever \textit{cualquiera} occurs under volitional verbs (see \citealt{ChoiRomero2008}, \citealt{am2018}). As I will show, neither approach accounts for the evaluative function of postnominal \textit{cualquiera}.

\citet{ChoiRomero2008} assume that \textit{cualquiera} has a similar denotation as \citegen{Fintel2000} denotation of \textit{wh-ever} in English. Von Fintel (\citeyear*{Fintel2000}) adopts \citegen{Dayal1998} insight that \textit{-ever} Free Relatives (FRs) introduce a layer of quantification over possible worlds and proposes that -\textit{ever} in -\textit{ever} FRs induces a presupposition of variation on either counterfactual worlds or epistemic worlds. Von Fintel (\citeyear*{Fintel2000}) points out that a subtype of \textit{-ever} FRs expresses “indifference” on somebody’s part. Compare \REF{bkm:Ref533762240} and \REF{ex:key:97b}. Both of them assert the same proposition paraphrasable using a definite description, namely, the proposition that the person who was at the top of the ballot won the election yesterday (see \citealt{ChoiRomero2008}: 88).

\largerpage
\ea  %ex 97
\ea \label{bkm:Ref533762240} In yesterday’s election, who was at the top of the ballot won.
\ex \label{ex:key:97b} In yesterday’s election, whoever was at the top of the ballot won.
\z
\z\clearpage

Unlike \REF{bkm:Ref533762240}, \REF{ex:key:97b} conveys an extra meaning triggered by \textit{-ever}, such that the identity of who was at the top of the ballot did not matter to winning yesterday’s election. Put differently, the identity of the denotation of \textit{-ever} FRs does not matter for the general nature or outcome of the election.

In \citet{Fintel2000}, this essential link follows from the presupposition of variation given in \REF{bkm:Ref533762335}, which is identified as the nature of \textit{-ever}’s contribution. The presupposition of variation tells us that if the individual denoted by an \textit{-ever} FR had been different, the truth value of the assertion in the actual world would still be valid in all the counterfactual worlds.

\ea %ex 98
\label{bkm:Ref533762335} Presupposition of variation for \REF{ex:key:97b}: If the person who was at the top of the ballot had been different, the same thing would have happened: that (new) person would have won.
\z

That is, the presupposition triggered by \textit{whoever was on the top of the ballot won} conveys that if the person who won the election had been different in all the counterfactual worlds, the truth of the proposition \textit{the person at the top of the ballot won} would also hold in the counterfactual worlds. From this presupposition of variation, it is inferred that regardless of who was at the top of the ballot, \textit{being at the top of the ballot} and \textit{winning yesterday’s election} are in an essential relation.

In \citet{Fintel2000}, a sentence containing an \textit{-ever} FR is formalised as in \REF{bkm:Ref533762404}. In the formulae, F indicates the modal base for \textit{-ever} FRs, which is a set of worlds on which the presupposition of variation operates. P refers to the denotation of the NP property contained in the \textit{-ever} FR, and Q refers to the property expressed by the rest of the sentence. Sentences containing an \textit{-ever} FR assert that the thing that has P is Q in the actual world, as shown in \REF{bkm:Ref533762404}. Note that, according to \citet{Fintel2000}, \textit{{}-ever} FRs denote definite descriptions on the assertion level. This is represented by the iota-operator ι. The presupposition triggered by \textit{-ever} says that in every world w' (of the corresponding modal base) that are different from the actual world only with respect to the referent of the \textit{-ever} FR, the asserted proposition has in w′ whatever truth value it has in the actual world w\textsubscript{0} (see \citealt{ChoiRomero2008}: 88):

\ea %ex 99
\label{bkm:Ref533762404}  Whatever (w\textsubscript{0})(F)(P)(Q) according to \citet{Fintel2000}
\ea  \textbf{Asserts}: Q(w\textsubscript{0})(ιx.P(w\textsubscript{0}) (x))
\ex \textbf{Presupposes}: ${\forall}$w′ ${\in}$ min\textsubscript{w0} [F ${\cap}$ λw″. ιx.P(w″)(x) ${\neq}$ ιx.P(w\textsubscript{0})(x)]: Q (w′) (ιx.P(w′)(x)) = Q (w\textsubscript{0}) (ιx.P(w\textsubscript{0})(x))
\z
\z

Choi \& Romero apply the formalisation in \REF{bkm:Ref533762404} to the sentence in \REF{bkm:Ref150779072}. The modal base F is counterfactual, and thus a presupposition of counterfactual variation is conveyed, as in \REF{ex:key:100b}.

\ea \label{bkm:Ref150779072} %ex  100
\ea In yesterday’s election, whoever was at the top of the ballot won.
\ea \textbf{Assertion} : λ w\textsubscript{0}. win(ι y.top-of-ballot(y,w\textsubscript{0}), w\textsubscript{0})
\ex \textbf{Presupposition} \label{ex:key:100b}: λ w\textsubscript{0} ${\forall}$ w′ ${\in}$ min\textsubscript{w0} [F ${\cap}$ λ w″. [ι y.top-of-ballot(y,w″ ${\neq}$ ι y.top-of-ballot(y,w\textsubscript{0})]: win (ι y.top-of-ballot(y,w′), w′) = win (ι y.top-of-ballot(y,w\textsubscript{0}), w\textsubscript{0})
\z
\ex In prose: 
\ea \textbf{Assertion}: In w\textsubscript{0}, the person who was at the top of the ballot in w\textsubscript{0} won.
\ex \textbf{Presupposition}: In each world w′, a counterfactual world of w\textsubscript{0}, if someone else had been at the top of the ballot in w′, the person who was at the top of the ballot in w′ won in w′ iff the person who was at the the top of the ballot in w\textsubscript{0} won in w\textsubscript{0}.
\z
\z
\z 

\citet{ChoiRomero2008} apply von Fintel’s formalisation of \textit{wh-ever} to the Spanish indefinite \textit{cualquiera} with a minor change: the ordinary meaning of Spanish indefinite \textit{cualquiera} contains an existential quantifier and not an iota operator as in von Fintel’s analysis of \textit{wh-ever}. The Spanish example in \REF{bkm:Ref517736} expresses the assertion that some book was picked and put on a pile and a presupposition of variation according to which there are other maximally similar worlds in which a book has been picked and put on a pile (see \citealt{ChoiRomero2008}: 91):

\ea \label{bkm:Ref517736} %ex 101 
\ea \gll Juan cogió un libro \textit{cualquiera} y lo puso encima de la pila. \\
Juan took a book \textsc{cualquiera} and it put on.top of the pile\\
\glt ‘Juan took some book and put it on top of the pile.’
\ea \textbf{Assertion}: λw\textsubscript{0}. ${\exists}$x.book(x,w\textsubscript{0}) \& pick(j,x,w\textsubscript{0}) \& put-on-pile(j,x,w\textsubscript{0})
\ex \textbf{Presupposition}: \\
λw\textsubscript{0} ${\forall}$ w′ ${\in}$ min\textsubscript{w0} [F ${\cap}$ λw″. \{x:book(x,w″)\} ${\neq}$ \{x: book(x,w\textsubscript{0})\}]:\\
${\exists}$x.book(x,w′) \& pick(j,x,w′) \& put-on-pile(j,x,w′)\\
${\exists}$x.book(x,w\textsubscript{0}) \& pick(j,x,w\textsubscript{0}) \& put-on-pile(j,x,w\textsubscript{0})
\z
\ex In prose:
\ea  \textbf{Assertion}: In the actual world w\textsubscript{0}, there is some book in w\textsubscript{0} that John picked up and put on the pile in w\textsubscript{0}.
\ex \textbf{Presupposition}: In all counterfactual worlds w′ minimally different from w\textsubscript{0} with respect to the identity of the set of books, there is some book in w′ that John picked up and put on the pile in w′ iff there is some book in w\textsubscript{0} that John picked up and put on the pile in w\textsubscript{0}.
\z
\z
\z

In contrast to \citet{ChoiRomero2008}, \citet{am2018} assume that the random choice interpretation introduces a modal domain consisting of worlds compatible with agent’s decisions. \citet{am2018} describe RC using a modal component presented in \REF{bkm:Ref150779122-2} according to which the relevant cards that Juan did not choose but could have chosen are part of the cards the agent had at the time he took the decision \textit{d\textsubscript{e}}. The modal component states that Juan’s decision could be fulfilled by taking any card in the domain. The following components are part of truth conditions (for the formalisation see \citealt{am2018}):

\ea % ex 102 
    \gll Juan cogió una carta \textit{cualquiera}.\\
    Juan took a card \textsc{cualquiera}\\
    \glt ‘Juan took a random card.’
\z


\ea \label{bkm:Ref150779122} %ex 103
\ea \textit{Existential Component}: In \textit{w}, there is an event \textit{e} of Juan taking a card \textit{x}.
\ex \label{bkm:Ref150779122-2} \textit{Modal Component}: For every (relevant) card \textit{y} in \textit{w}, there is a world \textit{w}′ compatible with \textit{d\textsubscript{e}} where there is an event \textit{e}′ of Juan taking \textit{y} that fulfills \textit{d\textsubscript{e}}.
    (\citealt{am2018}: 32)
\z
\z

I dub this approach \textit{intentional randomness}, as the agent who picked a random card had made the decision to do so.

\subsubsection{Discussion} \label{sec:2.2.3.3}

I would like to discuss three critical points concerning the previous analyses of \textit{cualquiera}. First, no study so far suggests a detailed analysis of the evaluative interpretation of \textsc{un} N \textit{cualquiera} in episodic and predicative contexts or makes a proposal as to how to analyse the relation between the evaluative meaning and RC/FC.

The second critical point concerns the description of the random choice interpretation. I will show in \chapref{sec:4} that \textit{cualquiera} can express unintentional/unvoluntary randomness too, which is the result of an unreflected/unplanned action. Let us consider the following scenario:

\ea \label{ex:key:104} %ex 104 
  Scenario: Imagine Juan is a lunatic. He sleepwalks and picks up several objects in his sleep, including some cards:

  \gll Juan cogió una carta \textit{cualquiera}.\\
  Juan took a letter \textsc{cualquiera}\\
  \glt ‘Juan took a random card.’
\z

In this scenario, \textit{cualquiera} has a RC interpretation even though the event of taking the actual card was not intentional and even though Juan did not decide to take a card. Moreover, I will show in \chapref{sec:4} on diatopic variation of \textit{cualquiera} in Argentinian Spanish, that there are volitional/intentional verbs that have an evaluative interpretation as their only meaning when they select \textit{cualquiera} as a direct object , such as \textit{mandar cualquiera} ‘to say bullshit’ (see also Fr. \textit{dire n’importe quoi} ‘to say bullshit’). The following examples, also from Argentinian Spanish, demonstrates such a case.

\ea \label{ex:key:105} %ex 105
  \gll Perdón respondí \textit{cualquiera}.\\
  sorry answer.1SG \textsc{cualquiera}\\
  \glt ‘Sorry, I answered \textit{nonsense}.’\\
  (@ 33, Buenos Aires/neighrborhood Congreso, 25 years old; Argentinian Spanish)
\z

\ea 
\label{ex:key:106} %ex 106
  \gll Osea respondí re \textit{cualquiera} lo que me preguntaba, no sé qué carajo, habré estado muy dormida pero como que contesté las preguntas por la mitad…\\
   well answered very \textsc{cualquiera} what that {me} asked \textsc{neg} know what hell must.have been very asleep but as if answered the questions by the half\\
  \glt ‘I answered with complete nonsense to what I was asked, I don’t know, what the hell, I must have been very tired because I answered the questions only partially.’

  (@ 34, female blogger, Buenos Aires, 25 years old; ArgSp)
\z

We thus need to find a better distinction between RCI and the evaluative interpretation that is based on more criteria than just the presence of a volitional/intentional verb. The relation between random choice and evaluative meaning needs to be explained. I address this issue in \chapref{sec:4}.

Another critical point of previous approaches to RCI is that they are too restrictive. Both approaches of RC in \sectref{sec:2.2} suggest that there is a correlation between volitional verbs and RC of postnominal \textit{cualquiera}. However, this is too restrictive. There are examples of \textsc{un} N \textit{cualquiera} in mathematical contexts under copular verbs or in subject position with free choice interpretation. Note that \textsc{un} N \textit{cualquiera} should not have RC interpretation in subject or predicative position under either approach in \sectref{sec:2.2}. Let us consider the following examples:

\ea %ex 107
\label{bkm:Ref43031271}  \gll El punto P de coordenadas (a,b) es un punto \textit{cualquiera}.\\
the point P of coordinates (a,b) is a point \textsc{cualquiera}\\
\glt ‘The point P of coordinates (a,b) is a random point.’

(@ 35)
\z

\ea %ex 108
\label{bkm:Ref43031168}  
\gll una muestra \textit{cualquiera} nos vale.\\
a sample \textsc{cualquiera} to.us is.worth\\ 
\glt ‘Any random sample is good for us.’

  (@ 36)
\z

In these examples, random choice is implicit; that is, what is random about the point or the sample is that the point P or the sample is defined in a way such that they were picked randomly or can be picked randomly. Random choice is thus not restricted to the occurrence of volitional verbs expressed in the sentence, but to contexts in which it makes sense to recover a choice. Note that in these examples, the evaluative reading of ‘ordinary, average, common’ does not make sense at all, because there is no point in the coordination system that can be described as unremarkable or common for the speaker in a mathematical context and it does not make sense to talk about an unremarkable sample in statistics. I thus disagree with \citet{Alonso-OvalleMenéndez-Benito2011,am2018}, who predict an unremarkable reading of \textsc{un} N \textit{cualquiera} in affirmative sentences in predicative position such as \REF{bkm:Ref43031271} and in subject position such as \REF{bkm:Ref43031168}.

\section{{\textit{Cualquiera}} {under} {negation}} \label{sec:2.3}

In this section, I will examine FCIs under negation more closely, as negation will play a role in the corpus data description of \textsc{un} N \textit{cualquiera} in \chapref{sec:3}.

The use of FCIs under negation was first described by \citet{Horn2000}. The speaker in \REF{bkm:Ref42805181} says that she wants to distinguish or discriminate among the possible candidates she wants to go to bed with. In other words, the speaker wants to be selective. According to Horn, \textit{not} negates the indiscriminacy. Thus, \textit{not just any} has an anti-indiscriminative interpretation (see \citealt{Horn2000}, \citealt{AloniPort2013}, among many others):

\ea \label{bkm:Ref42805181} %ex 109
 I don’t want to go to bed with just anyone anymore. I have to be attracted to them sexually.

  (\citealt{AloniEtAl2011}: 8)
\z

\citet{Fălăuş2013} offers a more general analysis of FCIs under negation. According to Fălăuş, Italian \textit{qualunque} in \REF{bkm:Ref42805226} has the denotation of an existential quantifier like \textit{some} in English (see Fălăuş 2013 for formal representation and \sectref{sec:2.1.3}). Moreover, \textit{qualunque} triggers alternatives \{a, b, c,…\}, such as alternative books \{\textit{Don Quijote}, \textit{War and Peace}, …\}. Postnominal \textit{qualunque} is focused, which triggers a domain restriction on the quantification of ‘some’. The sentence with a focused \textit{un libro qualunque} is interpreted as meaning that Gianni did not read a book in the initially considered domain of quantification (say D), but he read a book in some other domain (say D′), e.g. some exceptional book. The book in the domain D′ is singled out pragmatically, according to \citet{Fălăuş2013}.

\ea \label{bkm:Ref42805226}  %ex 110
\gll Gianni non ha letto un libro \textit{qualunque}.\\
  Gianni \textsc{neg} has read a book \textsc{qualunque}\\ 
  \glt ‘Gianni didn’t read just any book.’\\
  (Italian, \cite[95]{Falaus2013})
\z

The negation and focus interact in \REF{bkm:Ref42805226}. If \textit{qualunque} were not focused in \REF{bkm:Ref42805226}, negation would operate on {every} domain alternative and \textit{qualunque} would be interpreted as a negative polarity item similar to English \textit{any} under negation. It would get a different interpretation, namely the plain negation of \textit{Gianni did not read any book}. However, such a function of plain negation is expressed by negative words such as \textit{niente/nessun} ‘anything/any’ in Italian:

\ea \label{ex:key:111} %ex 111
  \gll Gianni non ha letto niente / nessun libro.\hfill (Italian)\\
Gianni \textsc{neg} has read nothing / no book\\
  \glt ‘Gianni did not ready anything/any book.’
\z

Instead, the negation in \REF{bkm:Ref42805226} operates on a {restricted} set of alternatives; that is, it negates every domain alternative out of a particular set of alternatives.

The same facts have been observed for Spanish \textit{cualquiera}, which is stressed under negation according to \citet{Arregui2006}, represented in capital letters as \textsc{cualquier} (see \citealt{Arregui2006}). If \textit{cualquier} is not stressed, it is ungrammatical, because it would be then interpreted as \textit{some} under negation as in \REF{bkm:Ref533953565}.

\ea \label{ex:kex:112} %ex 112
  \gll No compró \textit{cualquier} libro.\\
\textsc{neg} bought \textsc{cualquier} book\\
  \glt ‘S/he didn't buy just any book.’

  (\citealt{Arregui2006})
\z


\ea  \label{bkm:Ref533953565} %ex 113
  *No compró \textit{cualquier} libro.

  (intended NPI interpretation: ‘S/he didn’t buy any book.’)

  (\citealt{Arregui2006}: e.g. 2)
\z

There are other indefinite items that are used in the scope of negation with the NPI interpretation, namely negative indefinites such as \textit{ningún} or \textit{nadie}:

\ea \label{ex:key:114} %ex 114
 \gll  No vió a nadie.\\
\textsc{neg} saw to anyone\\
  \glt ‘S/he didn’t see anyone.’

  (@ 37)
\z

What remains unclear under \citet{Fălăuş2013} and \citegen{Arregui2006} proposal is the role of stress. It is a general rule that all FCIs, regardless of whether they are stressed or not and no matter what syntactic position and function they have, produce the anti-indiscriminative interpretation under negation (see \sectref{sec:2.2.2}). Stress is a side effect at best, not a necessary condition:

\ea 
\label{bkm:Ref42806294} %ex 115
\gll Mi padre no es \textit{cualquiera}.\\
my father \textsc{neg} is \textsc{cualquiera}\\
\glt ‘My dad is not just anyone.’\footnote{If stress played an important role, we would not observe the robust correlation between the anti-discriminative effect of FCIs under negation (see \sectref{sec:3}).} 

(@ 38)
\z


\ea \label{bkm:Ref42806303} %ex 116
  \gll Hoy no es \textit{cualquier} día, sino ese en el que todo puede cambiar.\\
today \textsc{neg} is \textsc{cualquier} day, but that in the which everything can change\\
  \glt ‘Today is not just any day, but that day when everything can change.’

  (@ 39)
\z

Another critical point is that this proposal does not say anything about the domain restriction of the indefinite \textit{qualunque}. What property do the books in the old domain D and new domain D′ in \REF{bkm:Ref42805226} have? How shall we account for the fact that the new domain D′ must contain a particular or a distinguished book that is not like any other books?

\ea \label{ex:key:117} %ex 117
\gll Gianni non ha letto un libro \textit{qualunque}. \# Ha letto un libro come tanti.  (Italian)\\
Gianni \textsc{neg} has read a book \textsc{qualunque} \# has read a book like many\\
  \glt ‘Gianni did not read just any book. \# He read a book like any other book.’
\z

Another approach to FCIs under negation is offered by \citet{Aloni2007}. According to her, FCIs under negation operate on the universal inference triggered by FCIs. The negation has scope over the conjunction of \textit{all} alternatives (it is not the case that you can bring along book A and you can bring along book B). This has the effect that the free choice is denied under negation, yielding the interpretation ‘you cannot freely choose to bring any book’:

\ea \label{bkm:Ref43030327} %ex 118
  \textit{No puedes traer \emph{cualquier} libro.}\\
  You cannot bring along just any book.

  ¬ [${\lozenge}$ you bring along book a1 or ${\lozenge}$ you bring along book a2 or … ], for any book a\textsubscript{i} ${\in}$ D

  ‘You cannot freely choose which book you bring me.’
\z


The result of the negation is the logical inference that you can bring along some book, but not all books. This logical inference is known in the literature as a \textit{quantitative scalar} \textit{implicature} (see \citealt{Chierchia2006,Chierchia2013}). To sum up, according to \citet{Aloni2007}, negation has scope over the universal quantification of free choice \textit{cualquiera} when \textit{cualquiera} is used under negation.

The prediction of this account is that FCIs under negation always produce the negation of free choice independent of stress. This is indeed the case, as pronominal \textit{cualquiera} produces the same effect (see \ref{bkm:Ref42806294} and \ref{bkm:Ref42806303}). As FC is denied under negation, free choice should not be interpreted as a conversational implicature, that is, it should only be available in certain pragmatic contexts, because it is not cancellable and does not disappear under negation in Modern Spanish. Since negation cannot cancel the free choice inference, this inference must then be either a presupposition or a conventional implicature (see \citealt{Potts2005} for tests of conventional implicature). Neither of them is cancellable. Both hypotheses exist in the literature: free choice is a presupposition (\citealt{Fintel2000}, \citealt{ChoiRomero2008}) or free choice is a conventional implicature (\citealt{Aloni2007,Aloni2021}). In this work, I ignore the status of the inference and focus instead on whether free choice is the \textit{only} lexical meaning of \textit{cualquiera} and whether other meanings can be derived from this lexical meaning.

Aloni’s account does not address about further implications of free choice alternatives under negation, such as the implication that the alternative candidates in \REF{bkm:Ref43030327} are interpreted as bad candidates on some pragmatic scale (see \citealt{Horn2000}). One could derive the following inference using free choice under negation in \REF{bkm:Ref43030327}. The negation of free choice entails careful or determined selection. The speaker in \REF{bkm:Ref43030327} thus asks to select a book carefully with determination. Careful selection of a book triggers the inference that the book in \REF{bkm:Ref43030327} must have the property of being distinct from other books that are not carefully selected (an interesting book, for instance). It is impossible to carefully select a book that does not stand out among other books. Imagine a scenario where all books are the same. In this scenario, it is impossible to select a book carefully. The hearer could thus interpret \textit{No puedes traer cualquier libro} in \REF{bkm:Ref43030327} as the negation of picking out a book with an undistinguished property and hence the imperative or advice to bring a distinctive book in \REF{bkm:Ref43030327}. As I will show, \textit{un libro} \textit{cualquiera} can denote a kind of a book, which has an undistinguished property of a book or a property of a book that all books share (see \chapref{sec:3}). As I will argue, it is this kind interpretation that has been lexicalised on \textsc{un} N \textit{cualquiera} in predicative contexts or episodic contexts such as \textit{es} \textsc{un} N \textit{cualquiera} (see \sectref{sec:3.3}).

\section{{Frequency} {of} {different} {interpretations} {of} {\textit{cualquiera}}} \label{sec:2.4}


\begin{figure}[b]
% \includegraphics[width=\textwidth]{figures/kellert-img001.png}
    \begin{tikzpicture}
        \begin{axis}[
            width  = \textwidth,
            height = .3\textheight,
            symbolic x coords = {CA,AA,CO,FC,UFC,GEN,INC,No-matter},
            axis lines*=left,
            bar width = 15pt,
            ybar,
            ymin=0,
            ymax=60,
            nodes near coords,
            ylabel={\%},
        ]
        \addplot+[lsRed] plot coordinates {
            (CA,3)
            (AA,1)
            (CO,9)
            (FC,58)
            (UFC,9)
            (GEN,5)
            (INC,8)
            (No-matter,9)
        };
        \end{axis}
    \end{tikzpicture}

\caption{\label{fig:key:2} Distribution of \textit{cualquiera} in Modern Spanish (i.e. 20th c. \citealt[19]{AloniEtAl2010})}

\end{figure}

In order to quantify different interpretations of Sp. \textit{cualquiera}, \citet{AloniEtAl2010} have classified corpus data from Corpus del Español (CdE) with respect to the classification system in \tabref{tab:key:fig}, which uses \citegen{Haspelmath1997} classification system of indefinites across languages that represent syntactic and semantic properties as demonstrated in \sectref{sec:2.1} %(\tabref{abc}).


\begin{table}[t]
\begin{tabularx}{\textwidth}{llllQ}
\lsptoprule
%  & \multicolumn{4}{c}{ {Functions on the map}}\\
&  &  {Abbr.} &  {Label} &  {Example}\\
\midrule
& a. & SK & specific known & \textit{Somebody} called. Guess who?\\
& b. & SU & specific unknown & I heard \textit{something}, but I couldn’t tell what it was.\\
& c. & IR & irrealis & You must try \textit{somewhere} else.\\
& d. & Q & question & Did \textit{anybody} tell you anything about it?\\
& e. & CA & conditional antecedent & If you see \textit{anything}, tell me immediately.\\
& f. & CO & comparative & In Freiburg, the weather is nicer than \textit{anywhere} else in Germany.\\
& g. & DN & direct negation & John didn’t see \textit{anybody}.\\
→ & h. & AM & anti-morphic & I don’t think that \textit{anybody} knows the answer.\\
→ & i. & AA & anti-additive & The gravity of such act goes beyond \textit{any} justification.\\
→ & j. & FC & free choice & \textit{Anybody} can solve this problem.\\
→ & k. & UFC & universal free choice & John kissed \textit{any} woman with red hair.\\
→ & l. & GEN & generic & \textit{Any} dog has four legs.\\
→ & m. & IND & indiscriminative & I do not want to go to bed with \textit{just anyone} anymore.\\
\lspbottomrule
\end{tabularx}
\caption{\label{tab:key:fig} Classification system of FCIs according to \citet{AloniEtAl2010}. Arrows mark contexts where any expressions appear with enriched interpretations, such as free choice, universal, or indiscriminative readings, often dependent on modal, generic, or negative operators.} %figure 1
\end{table}

As can be seen from the \figref{fig:key:2}, the most frequent function of all the possible contexts where \textit{cualquiera} appears in Modern Spanish in the 20th century is clearly the FC(I) function, that is, the distribution of total exhaustification of \textit{cualquiera} over alternatives (\figref{fig:key:2}). %is this the correct figure? 



The graph also shows that it can have other functions as well (e.g. the indiscriminative [IND] and ‘no matter’ interpretations). The latter interpretation is not defined in \tabref{tab:key:fig}. The meaning distribution misses the evaluative reading of \textit{cualquiera}, which I will investigate also with respect to its frequency in Modern Spanish in \chapref{sec:3}.

\section{{Summary} {of} {different} {properties} {of} {\textit{cualquiera}} {in} {synchrony}} \label{sec:2.5}

The critical review of the state of research in this chapter has shown that \textit{cualquier}a has different syntactic functions (determiner, nominal modifier, pronoun, noun, degree adjective), as well as different semantic functions, as summarised in \tabref{tab:key:4}.


\begin{table}
\footnotesize
\begin{tabularx}{\textwidth}{Q@{~}l@{~~}p{18mm}>{\raggedright}p{35mm}Q@{}}
\lsptoprule
{Semantic} {and} {syntactic} {properties} & {\textit{algún}} & {\textit{cualquier}} {NP{\slash}\newline \textit{cualquiera}} {as} {pronoun} & {NP} {\textit{cualquiera}} & {Source} \\
\midrule
{FCI/HI} & no & yes & Yes under modal verbs, subtrigging, comparatives, negation. & \citet{Falaus2013, Falaus2014}, \citet{Alonso-OvalleMenéndez-Benito2011,am2018} \citet{Aloni2007,Aloni2021}\\
{RCI} & no & no & Yes (in object position under agentive verbs)\newline No (in subject position, adverbs and/or under non-agentive verbs). & \citet{ChoiRomero2008}, \citet{Alonso-OvalleMenéndez-Benito2011,am2018}\\
\tablevspace
{IGN} & yes &  & Yes in Argentinian Spanish (like Italian \textit{qualsiasi}). & \citet{Chierchia2006}, \citet{Kellert-francia-2023}\\
\tablevspace
{EVAL} {1} {‘ordinary’} ‘bad’/{‘out} {of} {place’}\newline {EVAL} {2} & no & no in Sp.\newline  yes in ArgSp & Yes . for EVAL as ‘ordinary’\newline Distribution unrestricted in Sp., i.e. \textit{cualquiera} is lexically ambiguous and can appear under modal contexts as well \newline Yes in ArgSp for EVAL as ‘ordinary’ and ‘bad/out of place’\newline Distribution restricted in ArgSp under specific verbs such as copular verbs and verbs of saying. & \citet{Alonso-OvalleMenéndez-Benito2011,am2018}, \citet{Rivero2011,Fălăuş2015,Conde2011,Rizzosalierno2013}\\
\tablevspace
{NPI}  & yes & no & no & \citet{Arregui2006}, \citet{AloniEtAl2010}\\
\tablevspace
{Modifiable} {by} {adverbs} {like} {\textit{excepto /}}{\textit{casi}} & no & yes & ? & \citet{Arregui2006}\\
\tablevspace
{Sub-trigging}  & no & yes & ? & \citet{Quer2000}, \citet{AloniEtAl2010}\\
\lspbottomrule
\end{tabularx}

\caption{\label{tab:key:4} Semantic and syntacticproperties of \textit{cualquiera} in Modern Spanish according to previous studies}
\end{table} %Table 4e

When \textit{cualquiera} acts as a quantificational determiner or pronoun, it is used as a FCI with a restricted distribution (modal contexts, episodic contexts with subtrigging, epistemic modality, or negation) in European Spanish. It has a universal interpretation if embedded under a possibility modal. In non-modal environments, such as episodic contexts and predicative contexts in affirmative sentences, N \textit{cualquiera} has usually an evaluative interpretation of ‘unremarkable’ or a random choice interpretation in European Spanish. In predicative contexts in negative sentences, it expresses negation of free choice (the so-called ‘not just any’ interpretation). \textit{Cualquiera} has the status of a degree adjective in Argentinian Spanish.

\section{{Diachronic} {research} {on} {\textit{cualquiera}}} \label{sec:2.6}
\subsection{{Grammaticalisation} {of} {\textit{cualquiera}}}  \label{sec:2.6.1}

It is still a debate in the literature on the diachrony of \textsc{qualquier} and other indefinites in Romance as to whether they are an innovation of the Romance languages, in particular whether they have undergone a grammaticalisation path that led from separate words into a single compound (henceforth: the Innovation or Grammaticalisation Hypothesis) or whether they are directly related to Latin indefinites, either inherited or remodelled, i.e. by the substitution of elements in the course of language development or even as a conscious imitation of Latin models (henceforth: the Continuation Hypothesis; see \citealt{cp2009}, \citealt{Pato2012}, \citealt{Mackenzie2019}, \cite{Kellert-arias-2023}, and references therein). Under the Innovation Hypothesis, most previous diachronic studies concentrate on the question of how the relative pronoun \textit{qual} and the verb \textit{querer} ‘want, wish’ became grammaticalised into a single complex compound (see \citealt{Palomo1934}, \citealt{Lombard1938}–39/1947–48, \citealt{Meier1950}, \citealt{Rivero1988}, \citealt{Espinosa2006}, \citealt{cp2009}, among others). According to \citet{AloniEtAl2010}, the grammaticalisation process cannot be retraced in the diachronic corpus CdE because it took place prior to the written documentation of Old Spanish, or, in other words, the compound \textit{qualquier,-e,-a} already existed in the earliest periods in this corpus (see also \citealt{cp2009}, \citealt{Rivero1988}). I will review both hypotheses in this section.

\subsubsection{{Continuation hypothesis}} \label{sec:2.6.1.1}

According to some linguists, Romance FCIs are based on or and/or influnced by Latin FCIs, including calques from Latin (henceforth: the Continuation Hypothesis). As early as the late 19th century, \citet[57]{Meyerlubke1899} assumed that Italian \textit{qualunque} was remodelled on Latin \textit{qualiscumque} ‘(of) any (kind) whatever’, composed of \textit{qualis} ‘which’ + \textit{cumque} ‘ever’, where the second component was substituted with \textit{unquam} ‘any, any-’, hence \textit{qual[is] unqua[m]} (> Old and Modern Italian \textit{qualunque}) (see also \citealt{CortelazzoZolli1979}: 335; \citealt{Rohlfs1966-1968}, vol. II: 222–223). Etymological dictionaries and some linguists also assume that the French indefinite \textit{quiconque} has its source in Old French \textit{qui que} ‘whoever’ + \textit{onques} ‘ever’ (< Lat. \textit{unquam}), which was influenced by Latin \textit{quicumque} (see \citealt{Becker2014} citing \citealt{Godefroy2006}–1902, vol. 6: 511; \citealt{Gamillscheg1969}: 737; \citealt{Greimas1998}: 489; \citealt{Tobler1925-2002}–2018, vol. 8: 91). \citet{Menendezpidal1928} assumes that Old Spanish \textit{qual quier}, \textit{qui quier}, \textit{qual-se-quiera} represent the Old Spanish equivalents of the Latin indefinites \textit{quilibet, qualislibet}. \citet[387f.]{Meier1950} assumes that Old Sp. \textit{qualquiera} and \textit{quienquiera} ‘whoever’ were already fixed expressions in Old Spanish. However, he admits that the attested morphological variation of -\textit{quiera} (e.g. -\textit{quiere, -quisiere}) poses a problem for the fixed expression hypothesis. Another problem for the Continuation Hypothesis is that Old Spanish \textit{qualquier,-e,-a} contains a verb different from the verb inside the Latin indefinites, which have -\textit{libet} (lit. ‘pleases’),\footnote{{Latin
has also indefinite forms with the verb form} {\textit{vis}} {‘you want’:} {\textsc{quivis}}, {\textsc{quodvis}} {lit. ‘who you want’, ‘what you want’, etc. (see \citealt{Pato2012}: 304).}} with the consequence that, like in the case of It. \textit{qualunque} mentioned above, one has to assume a substitution with forms of (Old) Sp. \textit{querer} ‘to want, to love’. The Spanish verb \textit{querer} stems from Latin \textsc{quaerere} ‘to search, to investigate, try to find’ (\citealt{cp2009}), which has acquired the volitional meaning of ‘want’ (see \citealt[1152]{cp2009}, Fn.24). \citet[1152]{cp2009} assume that the free choice property can be more directly derived from ‘to try to find’ than from ‘to want, to desire’.\footnote{{This view is problematic for different reasons: First, in Catalan or Italian the Latin verb} {\textsc{volere}} {was chosen (}{\textit{qualsevol, qualsivoglia}}{), and, second, typologically, there is a whole range of volitional verbs that can feature in indefinites \citep{Haspelmath1997}, but to my knowledge, verbs meaning ‘to search’ are not among them.}} \citeauthor{cp2009} also note that Old Spanish \textit{qualquier,-e,-a} can be often found as a graphically separated form in the 13th century, a fact that they interpret as being against the Continuation Hypothesis (see \sectref{sec:5.1.2} for details).

\largerpage
  In most of the accounts that I have subsumed under the label Continuation Hypothesis, it remains somewhat unclear whether the authors believe in an organic development, in calques from written Latin, or in a combination of the two.

  \citegen{Pato2012} position on the origin of \textit{cualquiera} is not very clear. He seems to suggest that both the Continuation and the Innovation Hypotheses are correct in that — although \textit{qualquier,-e,-a} may be an innovation (but he does not explicitely say so) — Old Spanish inherited certain syntactic operations from Latin that can be applied to it. In particular, he is referring to the phenomenon of interposition, in which a noun or an NP is placed between the pronoun and the verb \textit{quier} (e.g. \textit{qual manera quier} lit. ‘which manner one wants’). He explicitly calls it “a structure inherited from Latin” \citep[273]{Pato2012}. Pato seems to believe that \textit{qual quier} was a quasi indefinite pronoun which could be separated by some linguistic material such as nouns, similar to ‘what manner soever’.

  However, Pato argues that the form without interposition, i.e. \textit{qual quier} N(P) ‘which one wants N(P)’ continued the path of grammaticalisation, becoming more fixed after the 13th century, as shown by graphically synthetic form \textit{qualquier} and later \textit{cualquier} \citep[297]{Pato2012}.

  I will first review Pato’s argumentation for the Continuation Hypothesis and then discuss the problems it entails.

  \citet{Pato2012} assumes that the structure with interposition is a characteristic phenomenon of Spanish in the 13th century, which had already arisen in Latin, where \textit{quilibet} “may be separated by tmesis into \textit{qui} and \textit{libet}” (\cite[62]{GildersleeveLodge1900}, quoted by \cite[304]{Pato2012}). The example given by Gildersleeve is Sallust, \textit{De coniuratione Catilinae} 5,4, which reads \textit{quoius rei lubet} ‘of any thing’.\footnote{The example itself is not quoted by either \citeauthor{GildersleeveLodge1900} or \citeauthor{Pato2012}, but was retrieved from the edition of Sallust’s works by \citet[5]{Ahlberg1919}. \textit{Quoius} is a variant of \textit{cuius}, genitive of the relative pronoun \textit{quī}, and \textit{lubet} is a variant of \textit{libet}. {Pato himself (\citeyear{Pato2012}: 305) tries to provide a further example from the late Latin}{} \textit{Lex Romana Visigothorum} (which was in use on the Iberian Peninsula): “[…] qui cuius libet rei possessione privati sunt […]”. However, this does not seem to be a case of interposition, if it means, as I think, “who [\textit{qui}] are deprived of the possession of any [\textit{cuius} \textit{libet}] thing [\textit{rei}].”} According to Pato, some Romance languages have inherited this structure. He suggests that the interposition corresponds to the classical poetic structure of tmesis, which is used to separate a syntactic structure for emphatic reasons: “la tmesis o \guillemotleft transposición de palabras\guillemotright{}
  (Institutio Oratoria, Quintiliano) es una licencia poética y retórica clásica (variante del hipérbaton) que consiste en la separación de los dos elementos que forman un compuesto”  ‘tmesis or “transposition of words” (Institutio Oratoria, Quintilian) is a classical and rhetorical license (a variant of hyperbaton) consisting of the separation of two elements forming a compound’ \citep[280]{Pato2012}.

  \citet{Pato2012} examines a corpus of works (directed) by Alfonso X of Castile (the Wise), where \textit{qual quier} mostly shows interposition of the following nouns: \textit{estrella} ‘star’, \textit{grado} ‘degree’, \textit{manera} ‘manner’, \textit{logar} ‘place’, and \textit{planeta} ‘planet’ (e.g. \textit{qual manera quier} ‘whichever manner’). The interposition occurs more often than the absence thereof (“sin interposición” ‘without interposition’), which he analyses as \textit{qual quier} NP (97.5\% vs 74\%, respectively). Pato observes that the form \textit{qual quier} also appears as a single word form in the 13th century. However, he only investigates the separated form  (i.e. \textit{qual quier manera} vs. \textit{qual manera quier}). The interposition seems mostly characteristic of the scientific texts of the corpus, in particular of the \textit{Los libros del saber de astronomía} and the \textit{Cánones de Albateni} \citep[301]{Pato2012}.


\begin{table}
\footnotesize
\begin{tabularx}{\textwidth}{lQrrrrrr}
\lsptoprule
%\hhline%%replace by cmidrule{~~------}
 &  & \multicolumn{2}{c}{  {Con}  {interposición}} & \multicolumn{2}{c}{  {Sin}  {interposición}} & \multicolumn{2}{c}{  {Totales}  {por}  {texto}}\\
 \midrule
 {1.} &  {\textit{Libros}  {del}  {saber}  {de}  {astronomía}} & 58\% & (334/578) & 19\% & (32/168) & 49\% & (366/746)\\
 {2.} &  {\textit{Cánones}  {de}  {Albateni}} & 15\% & (86/578) & {}-{}- & {}-{}- & 11.5\% & (86/746)\\
 {3.} &  {\textit{Espéculos}} & 10\% & (60/578) & 4.1\% & (7/168) & 9\% & (67/746)\\
 {4.} &  {{Formas}  {e}  {imágenes}} & 4.5\% & (26/578) & 9\% & (15/168) & 5.5\% & (41/746)\\
 {5.} &  {{Libro}  {del}  {cuadrante}  {señero}} & 4\% & (21/578) & {}-{}- & {}-{}- & 2.8\% & (21/746)\\
 {6.} &  {{Judizios}  {de}  {las}  {estrellas}} & 4\% & (21/578)) & 10\% & (17/168) & 5\% & (38/746)\\
 {7.} &  {{Lapidario}} & 2.4\% & (4/578) & 14\% & (24/168) & 5\% & (38/746)\\
 {8.} &  {{Libro}  {de}  {las}  {cruces}} & {}-{}- & (1/578) & 13\% & (22/168) & 3\% & (23/746)\\
 {9.} &  {{General}  {estoria,}  {I}} & {}-{}- & (1/578) & 4\% & (7/168) & 1\% & (8/746)\\
\midrule
\multicolumn{2}{c}{ {Total}  {number}  {per}  {structure}}  & 97.5\% & (564/578) & 74\% & (124/168) & 92\% & (688/746)\\
\lspbottomrule
\end{tabularx}
\caption{\label{tab:key:5} Distribution of relative clause structures with and without interposition across different scientific and didactic texts in Old Spanish. The table reports the percentage and raw frequency of each structure by text, as well as overall totals. Adapted from \citet[282]{Pato2012}.} %Table 5
\end{table}

Pato assumes that interposition is correlated not only with text genre (scientific texts) but also with a certain information-structural function in which the NP was focused in the position between \textit{qual} and \textit{quier}: “parece sensato pensar que la interposición podía emplearse para enfatizar el sustantivo temático, por tanto era una variación aprovechable, pero con carácter literario” ‘it seems sensible to think that the interposition could be used to emphasise the topic noun, therefore it was a usable variation, but with a literary character’ \citep[284]{Pato2012}.

There is a problem with Pato’s version of the {Continuation Hypothesis}. Apart from the fact that Pato is not explicit on the exact relation between the Romance and the supposed Latin structure, i.e. is it really a structure inherited from Latin (as he explicitely says) or an effect of Latinising writing, as other parts of his article may suggest? The fact that interposition occurs mostly in scientific texts would rather suggest the latter.\footnote{Alternatively, in particular considering the fact that both {\textit{Los libros del saber de astronomía}} {and} {\textit{Canones de Albateni}} {are based on Arabic sources, \citet[306]{Pato2012} does not exclude that the interposition phenomenon of} {\textit{qual quier}} {might be a calque from Arabic, as suggested by Fernández \citet{Fernandez1992}.}}

\subsubsection{{Innovation hypothesis/Grammaticalisation hypothesis}} \label{sec:2.6.1.2}

According to some linguists, wh-based FCIs such as Spanish \textit{cualquiera} developed from wh-relative clauses.

An early approach of this type is that by \citet{Lombard1938,Lombard1947}. Importantly, the author remarks that the verbal component of \textit{quienquiera} and similar formations (including our \textsc{qualquier}) cannot be a calque of Latin models, for the simple reason that \textit{quier(e)} and \textit{quiera} do not translate the corresponding verbal components of the Latin items. These are \textit{libet} ‘it pleases’ (with no logical subject) and \textit{vis} ‘you want’ (in 2nd person singular), whereas the Spanish formations show the verb \textit{querer} ‘to love, to want’ in the 3rd person (and in the subjunctive). The Spanish structure shows an indeterminate logical subject in the meaning of ‘one’ (a typically Romance usage, but not typical of Latin), which can sometimes be made more precise by using an impersonal reflexive construction (of the type \textit{quien se quier,-e,-a).} \citet{Lombard1938,Lombard1947} thus considers not only \textsc{qualquier} but also Sp. \textit{quienquiera, quien se quiera, dondequiera} ‘wherever’, as well as an impressive number of similar formations in many other Romance varieties, as pan-Romance innovations. According to \citet[187]{Lombard1938}, in the beginning, the verbal forms in these expressions were, in reality, relative clauses, where the choice of the verb was free (e.g. Old Sp. \textit{tenie gomito é lanzaba cual manjar que le daban}\footnote{Quoted from \citet[159]{Gessner1895}, where the example is identified as stemming from \textit{Libro de los exemplos (467 a m)}.} ‘He was sick, and thrust out whatever food they gave him’), but later the verbs meaning ‘to be’ and ‘to want’ yielded more or less fixed expressions, which could have a relative clause of their own, which could also become fixed (like in \textit{cualquiera que sea,} Port. \textit{como quer que seja}, Catalan, \textit{qualsevulla (qualsevol) que sìa,} It. \textit{qual vuolsi sia}\textbf{).}

\begin{quote}
Ainsi, l’indéfini relatif esp. \textit{cual que} a donné l’indéfini général \textit{cual [que] quiera}, qui a donné le nouvel indéfini relatif \textit{cualquiera que}, qui a donné le nouvel indéfini général \textit{cualquiera que sea}, qui peut même donner le nouvel indéfini relatif \textit{cualquiera que sea que.} De ces trois types d’indéfinis relatifs, seuls les deux premiers, \textit{cual que} et \textit{cualquiera que}, ont réellement une existence.

‘Thus, the relative indefinite esp. \textit{cual que} gave rise to the general indefinite \textit{cual [que] quiera}, which gave rise to the new relative indefinite \textit{cualquiera que}, which gave rise to the new general indefinite \textit{cualquiera que sea}, which may even give rise to the new relative indefinite \textit{cualquiera que sea que}. Of these three types of relative indefinite, only the first two, \textit{cual que} and \textit{cualquiera que}, actually exist.’

    (\citealt{Lombard1938}/39: 188)
\end{quote}

In a similar vein, more recent authors (\citealt{Company-CompanyPozas-Loyo2009}, \citealt{Haspelmath1997}, and others) claim that wh-based FCIs such as Sp. \textit{cualquiera} have emerged as a result of a grammaticalisation process in which free relative clauses were reanalysed as indefinite noun phrases (see also \citealt{Palomo1934}, \citealt{Company-CompanyPozas-Loyo2009}, \citealt{Alconchel2012}, \citealt{Brucart1999}: §7.5.7). However, there is no consensus on the precise process of grammaticalisation of Old Spanish \textit{qualquiera}. 

One possible explanation why there is no consensus is that the indefinite \textit{qualquiera} was already present in Old Spanish documents. It is thus difficult to confirm the hypothesis of the grammaticalisation on the basis of empirical proof.

The origin of Spanish \textit{cualquiera} according to some authors, is represented in \REF{ex:key:308} (cf. \citealt{Cuervo1893} in \citealt{cp2009}: 1086). Note that this grammaticalisation path is a simplified representation, which takes modern graphical representation as basis \citep{Elvira2007}.

%\begin{figure}

\ea 
\ea 
\gll Haga en él cual castigo quiera  \hfill{(Step 1)}\\
do.\textsc{imp} in him which punishment want-you\\ 
\glt ‘Do upon him the punishment which you may want.’

\ex Haga en él cual quiera castigo\\ 
$\rightarrow$ Free Relative with NP in situ (so-called Split-DP) \hfill (Step 2)

\ex Haga en él cualquier(a) castigo\\ 
$\rightarrow$ Postverbal Indefinite, monoclausal analysis \hfill (Step 3)\\
‘Do to him whatever punishment you want’

\z 
\z

%\caption{\label{fig:key:3} Grammaticalisation path of \textit{cualquier} in Spanish (\citealt{cp2009}: 1086)}
%\includegraphics[]{}
%\end{figure} %Figure 3

The most characteristic signs for the grammaticalisation of \textit{qualquier,-e,-a} according to \citet{cp2009} are briefly described as follows: The indefinite \textit{qual-quier,-e,-a} changed in frequency, in phonology (i.e. the loss of \textit{-e} in the form \textit{qual quier}), and in its morphosyntactic status (see also \citealt{Lehmann1995} for typical features of grammaticalisation). Its previous semantic interpretation was weakened as well; that is, its etymological interpretation has been lost and it acquired a new syntactic status. I will return to the frequency and especially to the semantic weakening in \sectref{sec:2.6.2}. As for the phonological weakening, the loss of \textit{e} in the form -\textit{quier} is not specific to the item at issue here but is due to a general phonological apocope that applied to all kinds of syntactic categories from the 13th to the 14th century. It can thus hardly be counted as evidence for grammaticalisation.

Given the list of arguments for the Romance \textit{Innovation Hypothesis}, which have been already convincingly discussed in \citet{Lombard1938,Lombard1947}, I assume that this hypothesis is stronger than the {Continuation Hypothesis}, i.e. Old Sp.\textbf{} \textit{qualquiera,-e,-a} is derived from a relative clause structure.

\subsection{{Semantic} {change} {in} {time}} \label{sec:2.6.2}

Not much research exists on the diachrony of the evaluative meaning of \textit{cualquiera} in European or Argentinian Spanish. According to \citet{Company2009}, the change in meaning of \textit{cualquiera} from FCI to EVAL is due to a \textit{subjectification process}, whereby the meaning of \textit{cualquiera} acquires a more subjective interpretation over time (see also \citealt{Traugott1989} for \textit{subjectification} as trigger of meaning change in English). \citet{Company2009} suggests that the free choice indefinite \textit{cualquiera} has undergone the change into the evaluative \textit{cualquiera}, triggering speaker-independent alternatives to a more subjective meaning dependent on the speaker’s evaluation. However, it remains to be seen what this subjective interpretation exactly means and how this change towards a more subjective meaning occurred. I will deal with this question in Chapters \ref{sec:4} and \ref{sec:5}. With respect to the diachrony of \textit{cualquiera} in European Spanish in \chapref{sec:5}, I will also study the grammaticalisation of \textit{cualquiera} in order to understand how the occurrence of the evaluative meaning fits into the grammaticalisation process, which according to the literature correlates with a loss and not with an enrichment of meanings of the grammaticalised construction (\citealt{Traugott1989}, \citealt{Heine1991}, among others). With respect to the diachrony of \textit{cualquiera} in Argentinian Spanish, \citet{Rizzosalierno2013} assumes that the potential linguistic precedent of the gradable predicate \textit{cualquiera} in Argentinian Spanish is possibly related to the gradable predicate \textit{cualunque}, an Italianism (see \sectref{sec:2.1.4}). \textit{Cualunque} appeared in Argentinian Spanish during the Italian immigration (on \textit{qualunque} in Italian, see \citealt{Chierchia2006}, \citealt{AloniEtAl2010}, among others) (see \sectref{sec:2.1.4}). Like \textit{cualquiera} in Argentinian Spanish, \textit{cualunque} also allows degree modification (unlike in Italian, see \cite{Kellert-francia-2023}) and can be used predicatively or attributively in a postnominal position, as happens with \textit{cualquiera}, and have an evaluative interpretation like ‘ordinary’:

\ea \label{ex:key:119} %ex 119
  \gll El que yo digo \textit{era} \textit{re} \textit{cualunque.}\\
  the that I say was very \textsc{cualunque}\\
  \glt ‘The one that I said was very ordinary.’\\
  (\citealt{Rizzosalierno2013}: 28)
\z 

\ea \label{ex:key:120} %ex 120
  \gll Es una flaca \textit{cualunque} que se cree que es Marilyn Monroe.\\
  is a woman \textsc{cualunque} that \textsc{refl} believes that is Marilyn Monroe\\
  \glt ‘She is an ordinary woman who thinks she’s Marilyn Monroe.’\\
  (\citealt{Rizzosalierno2013}: 29)
\z

However, the change from an indefinite pronoun/determiner/postnominal modifier \textit{cualunque} to a gradable adjective \textit{cualunque} and for that matter the transfer of this function to \textit{cualquiera} remains unexplained. I will deal with this change in \chapref{sec:4}.

With respect to the question how different meanings developed, it remains an open question how the reanalysis of the indefinite \textit{cualquiera} into an item with the evaluative meaning ‘common/unremarkable’ in European Spanish and into a degree adjective with meaning ‘unremarkable/bad’ or an evaluative noun with the meaning ‘nonsense’ in Argentinian Spanish fits models of linguistic reanalysis in the frameworks of formal semantics (\citealt{Eckardt2006}, \citealt{Bar-asher-siegal2020}, \citealt{Aloni2021}, among others). According to these authors, it is possible to describe reanalysis on the level of morphosyntactic form (F) and/or on the level of meaning (M) as shown in \REF{bkm:Ref65487536}.

% With respect to the question how different meanings developed, it remains an open question as to how the reanalysis of the indefinite \textit{cualquiera} into an item with the evaluative meaning ‘common/unremarkable’ in European Spanish and into a degree adjective with meaning ‘unremarkable/bad’ or an evaluative noun with the meaning ‘nonsense’ in Argentinian Spanish fits models of linguistic reanalysis in the frameworks of formal semantics (\citealt{Eckardt2006}, \citealt{Bar-asher-siegal2020}, \citealt{Aloni2021}, among others). According to these authors, it is possible to describe reanalysis on the level of morphosyntactic form (F) and/or on the level of meaning (M) as shown in \REF{bkm:Ref65487536}.

\ea \label{bkm:Ref65487536}  %ex 121
Scenarios of reanalysis (see \citealt{Aloni2021}):

  Scenario 1: \{F1;M1\}\textsuperscript{t1} $\rightarrow$  \{F2;M2\}\textsuperscript{t2}

  Scenario 2: \{F1;M1\}\textsuperscript{t1} $\rightarrow$  \{F1;M2\}\textsuperscript{t2}

  Scenario 3: \{F1;M1\}\textsuperscript{t1} $\rightarrow$  
  \{F2;M1\}\textsuperscript{t2}
\z

Reanalysis can take place both at the morphological/syntactical level F and at the semantic level M (Scenario 1), or it can be restricted to one level (Scenarios 2 and 3). Diachronic semantics deals with changes in Scenarios 1 and 2. The important claim of the formal semantic studies is that, at both points of time (t1 and t2), F suits M in a compositional manner — this is necessary for a reanalysis to take place (see \citealt{Eckardt2006}, \citealt{Bar-asher-siegal2020}, \citealt{Aloni2021}).

\section{{Open} {questions} {on} {\textit{cualquiera}}{} {and} {outlook} {for} {this} {book}} \label{sec:2.7}

 In this section, I discuss open questions and present the outlook for this book.

\subsection{{Questions} {on} {synchrony}} \label{sec:2.7.1}

There is no detailed discussion or definition of the non-modal meaning of \textit{cualquiera} ‘ordinary/common’ (which I have called the evaluative meaning, EVAL1) in predicative contexts in Modern Spanish and how this meaning of \textit{cualquiera} is related to other meanings such as the free choice meaning ‘any’ and the random choice meaning ‘random’. The examples discussed in the literature seem to indicate that this evaluative meaning of \textit{cualquiera} depends on the verb class, i.e. it is specifically found with non-volitional verbs (e.g. \textit{copular} verbs, episodic tense, perception verbs) (see \sectref{sec:2.3}). EVAL2 is also present with verbs such as \textit{hacer/decir cualquiera} ‘do/talk nonsense’ in Argentinian Spanish and French \textit{n’importe quoi} (see \sectref{sec:2.1}). I will therefore look into verb classes more systematically and address the question how the definition of EVAL can capture the variation in meanings under predicative verbs such as the meaning ‘common’ and the deragotary version of it, such as ‘nothing special/unremarkable’ (EVAL1) and the meaning as ‘nonsense’ under volitional verbs such as \textit{hacer/decir} (EVAL2). This meaning variation correlates to the variation in varieties of Spanish (Argentinian vs. European Spanish), as the meaning ‘nonsense’ (EVAL2) is only present in Argentinian Spanish (see \chapref{sec:3} for European Spanish and \chapref{sec:4} for Argentinian Spanish).

Another question that needs to be answered is what syntactic and semantic status the postnominal \textit{cualquiera} has in European Spanish. Is it a determiner or something else? And what status does the evaluative \textit{cualquiera} have under volitional verbs in Argentinian Spanish? Is it an indefinite pronoun or a noun?

It also remains to be seen whether all readings of \textit{cualquiera} or at least a subset of them can be derived from a more general meaning of \textit{cualquiera} or whether \textit{cualquiera} is lexically ambiguous as \citet{am2018} suggest (see \sectref{sec:2.2.3}). My working hypothesis is that the evaluative meaning is the result of a reanalysis of the {\textit{modal}} \textit{cualquiera} (i.e. \textit{cualquiera} that is used under a modal) in certain syntactic and semantic contexts in diachrony. I will address these issues within the synchronic and diachronic study of \textit{cualquiera} in Chapters \ref{sec:3}, \ref{sec:4} and \ref{sec:5}.

The classification system of \textit{cualquiera} chosen by \citet{AloniEtAl2010} in \sectref{sec:2.5} is not satisfactory, because it mixes up syntactic and semantic properties of \textit{cualquiera} and does not correlate them. For instance, conditional antecedents and comparative clauses are syntactic structures that can license \textit{cualquiera}. Consequently, conditional antecedents and comparative clauses are syntactic contexts, as both sentence types can license FCIs. From the results in \figref{fig:key:2} in \sectref{sec:2.5}, we do not learn whether different functions correlate with different syntactic positions of \textit{cualquiera} (e.g. +/– prenominal position). The classification map in \tabref{tab:key:fig} in \sectref{sec:2.5} does not contain a definition of the ‘no matter’ interpretation, making it difficult to judge how it is different from the indiscriminative interpretation (IND). Moreover, the RC and evaluative interpretations are absent from the classification system as well as from the frequency distribution in \tabref{tab:key:fig} and \figref{fig:key:2} in \sectref{sec:2.1} and \sectref{sec:2.2.3.3}. Given the critical points mentioned with respect to Table 2.5, I decided to use a more detailed syntactic classification of \textit{cualquiera} for my corpus analysis of the word in synchrony and diachrony (see Chapters \ref{sec:3}, \ref{sec:4}, and \ref{sec:5}). I will describe the evaluative function of \textit{cualquiera} in more detail in Chapters \ref{sec:3} and \ref{sec:4} and relate it to its grammatical function (i.e. N \textit{cualquiera} in \chapref{sec:3} and \textit{decir cualquiera} in \chapref{sec:4}). The diachronic study in Chapters \ref{sec:4} and \ref{sec:5} will help us understand the syntactic and semantic status of \textit{cualquiera} in different positions and different varieties of Spanish.

\subsection{{Questions} {on} {diachrony} {and} {the} {diatopic} {variation} {of} {\textit{cualquiera}}} \label{sec:2.7.2}

One important diachronic question is how the evaluative meanings \textit{cualquiera} came into existence in European and Argentinian Spanish. Does it have the same origin? In order to answer this question I will study the diachrony of \textit{cualquiera} with EVAL2 in \chapref{sec:4} and the diachrony of EVAL1 in \chapref{sec:5}. According to many linguists (\citealt{Blank1997}, \citealt{Traugott1989}, \citealt{Eckardt2006}, among others), semantic change is often triggered by pragmatic enrichment of the target utterance. For instance, the new meaning can be the result of reanalysing the old meaning with contextual information (some pragmatic implicatures). The open question is whether the evaluative meaning of \textit{cualquiera} in European and Argentinian Spanish is based on some pragmatic reanalysis in a compositional manner.

I will argue in Chapters \ref{sec:3}, \ref{sec:4}, and \ref{sec:5} that the evaluative meaning of \textit{cualquiera} is the result of using free choice indefinites in certain syntactic positions, such as the predicative position. I will show that the appearance of the evaluative meaning of \textit{cualquiera} is an example of Scenario 1 in \REF{bkm:Ref65487536}. Moreover, I will argue that the evaluative meaning is the result of pragmatic strengthening because the free choice meaning (the original meaning) is enriched by contextual information in certain pragmatic contexts in which some attitude holder x expresses her evaluation towards the proposition ‘every possibility is an option’ (FCI) or towards the proposition ‘every possibility matches agent’s decisions’ (RCI) expressed by \textit{cualquiera}. This contextual interpretation has triggered a reanalysis of \textit{cualquiera} into an evaluative predicate in European and Argentinian Spanish with further reanalysis into a degree predicate in Argentian Spanish. I dub my analysis ‘modal indefinites under attitudes’. My analysis thus combines research from two different camps: modal and polarity sensitive items as discussed in this section and semantics dealing with attitudes and evaluation (e.g. personal taste predicates) (see \citealt{Heim1982}, \citealt{Lasersohn2005}, among others), which I will introduce in Chapters \ref{sec:3} and \ref{sec:4}.
